
% !TEX program = lualatex
% 巨大LuaLaTeXパッケージ総合ストレステスト
% 目的: 日本語 + 数式 + 図 + 表 + コード + 箇条書き + 化学 + 電気回路 + 索引 + 相互参照を一つの長大TeXに詰め込む。
% コンパイル: lualatex huge_lualatex_package_stress_test.tex を2回推奨。
\documentclass[12pt,a4paper,openany]{ltjsbook}

% ===== LuaLaTeX / Japanese core =====
\usepackage{luatexja}
\usepackage{luatexja-fontspec}
\usepackage{fontspec}

% ===== Page layout / typography =====
\usepackage[margin=22mm,headheight=20pt,footskip=12mm]{geometry}
\usepackage{setspace}
\usepackage{indentfirst}
\usepackage{fancyhdr}
\usepackage{lastpage}
\usepackage{titling}
\usepackage{titlesec}
\usepackage{tocloft}
\usepackage{enumitem}
\usepackage{multicol}
\usepackage{paracol}
\usepackage{needspace}
\usepackage{afterpage}
\usepackage{pdflscape}

% ===== Math =====
\usepackage{amsmath,amssymb,amsfonts,amsthm,mathtools,bm,bbm}
\usepackage{empheq}
\usepackage{cancel}
\usepackage{cases}
\usepackage{braket}
\usepackage{tensor}
\usepackage{esint}

% ===== Graphics / drawing =====
\usepackage[svgnames,dvipsnames,table]{xcolor}
\usepackage{graphicx}
\usepackage{float}
\usepackage{wrapfig}
\usepackage{subcaption}
\usepackage{caption}
\usepackage{tikz}
\usetikzlibrary{arrows.meta,calc,positioning,shapes.geometric,shapes.misc,decorations.pathmorphing,decorations.pathreplacing,patterns,fit,matrix,graphs,quotes,angles,intersections,backgrounds,3d,automata,petri,chains,mindmap,trees}
\usepackage{pgfplots}
\pgfplotsset{compat=1.18}
\usepackage[europeanresistors,americaninductors]{circuitikz}

% ===== Tables =====
\usepackage{booktabs,longtable,tabularx,array,multirow,makecell,colortbl,dcolumn,diagbox}
\usepackage{rotating}
\usepackage{threeparttable}
\usepackage{threeparttablex}

% ===== Boxes / framed content =====
\usepackage[most]{tcolorbox}
\usepackage{mdframed}
\usepackage{fancybox}
\usepackage{framed}

% ===== Code / algorithms =====
\usepackage{listings}
\usepackage{fancyvrb}
\usepackage[ruled,vlined,linesnumbered]{algorithm2e}

% ===== Science / units / chemistry =====
\usepackage{siunitx}
\sisetup{detect-all,per-mode=symbol,output-decimal-marker={.}}
\usepackage[version=4]{mhchem}
\usepackage{chemfig}

% ===== Utilities =====
\usepackage{xparse}
\usepackage{ifthen}
\usepackage{etoolbox}
\usepackage{xspace}
\usepackage{makeidx}
\usepackage{marginnote}
\usepackage{datetime2}
\usepackage{lipsum}
\usepackage{blindtext}
\usepackage{csquotes}

% ===== Hyperlinks last-ish =====
\usepackage{hyperref}
\usepackage{bookmark}
\usepackage[nameinlink,noabbrev]{cleveref}
\hypersetup{
  unicode=true,
  colorlinks=true,
  linkcolor=MidnightBlue,
  citecolor=ForestGreen,
  urlcolor=DarkRed,
  pdftitle={Huge LuaLaTeX Package Stress Test},
  pdfauthor={Generated by ChatGPT},
  pdfsubject={LuaLaTeX preamble and package coverage stress test}
}

% ===== Index =====
\makeindex

% ===== Page style =====
\pagestyle{fancy}
\fancyhf{}
\lhead{Huge LuaLaTeX Stress Test}
\rhead{\leftmark}
\cfoot{\thepage/\pageref{LastPage}}

% ===== Heading style =====
\titleformat{\chapter}[display]{\Huge\bfseries\sffamily}{\filleft\chaptertitlename~\thechapter}{1ex}{\titlerule\vspace{1ex}\filright}[
\vspace{1ex}\titlerule]
\titleformat{\section}{\Large\bfseries\sffamily}{\thesection}{1em}{}
\titleformat{\subsection}{\large\bfseries\sffamily}{\thesubsection}{1em}{}

% ===== Colors =====
\definecolor{mainblue}{HTML}{1D4ED8}
\definecolor{softblue}{HTML}{EFF6FF}
\definecolor{softgreen}{HTML}{ECFDF5}
\definecolor{softred}{HTML}{FEF2F2}
\definecolor{ink}{HTML}{111827}
\definecolor{codebg}{HTML}{F8FAFC}

% ===== Theorems =====
\theoremstyle{definition}
\newtheorem{definition}{定義}[chapter]
\newtheorem{example}{例}[chapter]
\theoremstyle{plain}
\newtheorem{theorem}{定理}[chapter]
\newtheorem{lemma}{補題}[chapter]
\theoremstyle{remark}
\newtheorem{remark}{注意}[chapter]

% ===== tcolorbox styles =====
\tcbset{
  enhanced,
  sharp corners,
  boxrule=0.6pt,
  coltitle=black,
  fonttitle=\bfseries,
  before skip=8pt,
  after skip=8pt,
}
\newtcolorbox{bluebox}[1]{colback=softblue,colframe=mainblue,title={#1},breakable}
\newtcolorbox{greenbox}[1]{colback=softgreen,colframe=ForestGreen,title={#1},breakable}
\newtcolorbox{redbox}[1]{colback=softred,colframe=FireBrick,title={#1},breakable}
\newtcbtheorem[number within=chapter]{boxedtheorem}{箱入り定理}{colback=softblue,colframe=mainblue,fonttitle=\bfseries,breakable}{bt}

% ===== Listings =====
\lstdefinestyle{pythonstyle}{
  language=Python,
  backgroundcolor=\color{codebg},
  basicstyle=\ttfamily\small,
  keywordstyle=\bfseries\color{MidnightBlue},
  commentstyle=\itshape\color{ForestGreen},
  stringstyle=\color{DarkRed},
  numbers=left,
  numberstyle=\tiny\color{Gray},
  stepnumber=1,
  numbersep=8pt,
  breaklines=true,
  frame=single,
  rulecolor=\color{Gray},
  tabsize=2,
  showstringspaces=false
}
\lstdefinestyle{texstyle}{
  language=[LaTeX]TeX,
  backgroundcolor=\color{codebg},
  basicstyle=\ttfamily\small,
  keywordstyle=\bfseries\color{MidnightBlue},
  commentstyle=\itshape\color{ForestGreen},
  numbers=left,
  numberstyle=\tiny\color{Gray},
  breaklines=true,
  frame=single,
  showstringspaces=false
}

% ===== Custom commands =====
\newcommand{\TeXDOM}{\textsf{TeX-DOM}\xspace}
\newcommand{\dd}{\,\mathrm{d}}
\newcommand{\ii}{\mathrm{i}}
\newcommand{\ee}{\mathrm{e}}
\newcommand{\R}{\mathbb{R}}
\newcommand{\C}{\mathbb{C}}
\newcommand{\N}{\mathbb{N}}
\newcommand{\Z}{\mathbb{Z}}
\newcommand{\vect}[1]{\bm{#1}}
\newcommand{\mat}[1]{\mathbf{#1}}
\DeclareMathOperator{\Tr}{Tr}
\DeclareMathOperator{\rank}{rank}
\DeclareMathOperator{\Var}{Var}
\DeclareMathOperator{\Cov}{Cov}
\DeclareMathOperator{\diag}{diag}
\DeclarePairedDelimiter{\abs}{\lvert}{\rvert}
\DeclarePairedDelimiter{\norm}{\lVert}{\rVert}
\DeclarePairedDelimiter{\inner}{\langle}{\rangle}

\title{巨大LuaLaTeXパッケージ総合ストレステスト}
\author{Generated for TeX64 / Fermion-style engine testing}
\date{\today}

\begin{document}
\frontmatter
\maketitle

\begin{abstract}
これは、LuaLaTeXでコンパイル可能な巨大TeXのテスト文書である。日本語、英語、数式、表、長表、TikZ、PGFPlots、CircuitTikZ、tcolorbox、mdframed、コード表示、アルゴリズム、化学式、単位、索引、相互参照、横向きページ、段組、脚注、余白注などを一つのファイルに詰め込んでいる。目的は、エディタ、プレビューエンジン、SyncTeX、補完、構文解析、アウトライン生成、折りたたみ、差分検出の総合負荷試験である。
\end{abstract}

\tableofcontents
\listoffigures
\listoftables

\chapter{プリアンブルの利用パッケージ一覧}
\begin{multicols}{2}
\begin{itemize}[nosep]
  \item 日本語: luatexja, luatexja-fontspec
  \item 文字: fontspec
  \item レイアウト: geometry, fancyhdr, titlesec, tocloft
  \item 数式: amsmath, mathtools, amsthm, bm, braket, tensor, esint
  \item 図: graphicx, tikz, pgfplots, circuitikz
  \item 表: booktabs, longtable, tabularx, multirow, makecell, diagbox
  \item 箱: tcolorbox, mdframed, fancybox, framed
  \item コード: listings, fancyvrb, algorithm2e
  \item 科学: siunitx, mhchem, chemfig
  \item 参照: hyperref, bookmark, cleveref
  \item その他: makeidx, lipsum, blindtext, csquotes, paracol, pdflscape
\end{itemize}
\end{multicols}

\mainmatter
\chapter{巨大総合章 1: パッケージ横断テスト}\label{chap:1}


\begin{redbox}{章の目的}
この章は、同じパターンを少しずつ変えながら反復し、アウトライン、ラベル、図表番号、ページ分割、脚注、索引、長い数式、コード、図形描画を大量に発生させるための章である。
\end{redbox}


\section{総合テストセクション 1-1}\label{sec:ch1-s1}
この節では、\TeXDOM、LuaLaTeX、数式組版、図表、相互参照、索引生成を同時に試験する。索引語としてLuaLaTeX\index{LuaLaTeX}、TikZ\index{TikZ}、tcolorbox\index{tcolorbox}、数式\index{数式}、日本語組版\index{日本語組版}を登録する。

\begin{definition}[状態ベクトル]
複素ヒルベルト空間 $\mathcal{H}$ の単位ベクトル $\ket{\psi}\in\mathcal{H}$ を純粋状態という。射影 $\rho=\ket{\psi}\bra{\psi}$ は $\rho^2=\rho$, $\Tr\rho=1$ を満たす。
\end{definition}

\begin{theorem}[スペクトル分解の簡易形]
任意のエルミート行列 $A\in\C^{n\times n}$ には、ユニタリ行列 $U$ と実対角行列 $D=\diag(\lambda_1,\dots,\lambda_n)$ が存在して
\begin{equation}\label{eq:spectral-1-1}
A=UDU^\dagger,\qquad \ee^{-\ii t A}=U\diag(\ee^{-\ii t\lambda_1},\dots,\ee^{-\ii t\lambda_n})U^\dagger
\end{equation}
と書ける。
\end{theorem}

\begin{proof}
有限次元のスペクトル定理による。相異なる固有値に属する固有空間は直交し、各固有空間で正規直交基底を選ぶことで $U$ が得られる。\cref{eq:spectral-1-1} は関数計算の定義から従う。
\end{proof}

\begin{empheq}[box=\fbox]{align}
Z(\beta)&=\Tr \ee^{-\beta H}=\sum_n \ee^{-\beta E_n},\\
F(\beta)&=-\frac{1}{\beta}\log Z(\beta),\\
\langle E\rangle&=-\frac{\partial}{\partial\beta}\log Z(\beta).
\end{empheq}

\begin{bluebox}{日本語・英語・記号混在テスト}
これは長い日本語段落である。The quick brown fox jumps over the lazy dog. 数式 $\int_{-\infty}^{\infty}\ee^{-x^2}\dd x=\sqrt{\pi}$ と、化学式 \ce{2H2 + O2 -> 2H2O} と、単位 \SI{6.62607015e-34}{\joule\second} を同時に含む。さらに、\enquote{引用符}、脚注\footnote{脚注テスト。日本語・English・$E=mc^2$ を含む。}、余白注\marginnote{Margin note 1-1}も入れる。
\end{bluebox}

\subsection{行列・テンソル・微分形式}
\begin{align}
\mat{M}_{1} &=
\begin{pmatrix}
1 & 1 & 1 \\
0 & 1 & 1+1 \\
1 & 0 & 1
\end{pmatrix},
&
\det \mat{M}_{1}&=1-1(1)(1+1),\\
\tensor{R}{^\rho_{\sigma\mu\nu}}&=\partial_\mu \Gamma^\rho_{\nu\sigma}-\partial_\nu\Gamma^\rho_{\mu\sigma}+\Gamma^\rho_{\mu\lambda}\Gamma^\lambda_{\nu\sigma}-\Gamma^\rho_{\nu\lambda}\Gamma^\lambda_{\mu\sigma},\\
\oint_{\partial\Omega}\vect{F}\cdot\vect{n}\dd S&=\iiint_{\Omega}\nabla\cdot\vect{F}\dd V.
\end{align}

\subsection{表と長いセル}
\begin{table}[H]
\centering
\caption{booktabs, makecell, multirow, colortbl の混在テスト 1-1}\label{tab:mix-1-1}
\begin{tabular}{llcr}
\toprule
\rowcolor{softblue} 分類 & パッケージ & 典型用途 & 負荷 \\
\midrule
数式 & amsmath/mathtools & \makecell{多行整列\\括弧制御} & 高 \\
図 & TikZ/PGFPlots & ベクター図形 & 高 \\
表 & longtable/tabularx & ページ分割表 & 中 \\
コード & listings/fancyvrb & ソース表示 & 中 \\
\bottomrule
\end{tabular}
\end{table}

\subsection{TikZ図}
\begin{figure}[H]
\centering
\begin{tikzpicture}[node distance=18mm,>=Latex]
\node[draw,rounded rectangle,fill=softblue,minimum width=26mm,minimum height=10mm] (src) {Source TeX};
\node[draw,rectangle,fill=softgreen,right=of src,minimum width=26mm,minimum height=10mm] (parse) {Parser};
\node[draw,diamond,aspect=2,fill=softred,right=of parse,minimum width=28mm] (dom) {TeX-DOM};
\node[draw,rounded rectangle,fill=softblue,right=of dom,minimum width=26mm,minimum height=10mm] (pdf) {PDF};
\draw[->,thick] (src) -- node[above]{tokens} (parse);
\draw[->,thick] (parse) -- node[above]{nodes} (dom);
\draw[->,thick] (dom) -- node[above]{layout} (pdf);
\draw[->,bend left=35,dashed] (pdf.north) to node[above]{SyncTeX} (src.north);
\end{tikzpicture}
\caption{TikZによる処理系模式図 1-1}\label{fig:tikz-1-1}
\end{figure}

\subsection{PGFPlots}
\begin{figure}[H]
\centering
\begin{tikzpicture}
\begin{axis}[width=.8\linewidth,height=55mm,grid=major,xlabel={$x$},ylabel={$f(x)$},legend pos=north east]
\addplot[domain=0:6.28,samples=80] {sin(deg(x))};
\addlegendentry{$\sin x$}
\addplot[domain=0:6.28,samples=80,dashed] {cos(deg(x))};
\addlegendentry{$\cos x$}
\addplot[domain=0:6.28,samples=80,dotted] {exp(-x/4)*sin(deg(3*x))};
\addlegendentry{$e^{-x/4}\sin(3x)$}
\end{axis}
\end{tikzpicture}
\caption{PGFPlotsテスト 1-1}\label{fig:pgfplots-1-1}
\end{figure}

\subsection{コードとアルゴリズム}
\begin{lstlisting}[style=pythonstyle,caption={Python風擬似コード 1-1}]
from math import exp

def partition_function(beta, energies):
    return sum(exp(-beta * e) for e in energies)

def free_energy(beta, energies):
    z = partition_function(beta, energies)
    return -log(z) / beta
\end{lstlisting}

\begin{algorithm}[H]
\caption{差分プレビュー更新アルゴリズム 1-1}
\KwIn{old DOM, new source tokens}
\KwOut{minimal changed regions}
Parse changed token interval\;
Update syntax tree nodes\;
Compute layout invalidation frontier\;
\If{math box changed}{Rebuild equation box and update cross references\;}
\Else{Reuse previous node metrics\;}
Return changed PDF rectangles\;
\end{algorithm}

\subsection{化学・回路・箱}
\begin{greenbox}{化学式と単位}
水の生成反応は \ce{2H2 + O2 -> 2H2O} であり、標準状態の議論では温度 \SI{298.15}{\kelvin}、圧力 \SI{1}{\bar} を使う。構造式の簡易表示: \chemfig{H-O-H}。
\end{greenbox}

\begin{figure}[H]
\centering
\begin{circuitikz}[scale=0.9]
\draw (0,0) to[V,l=$V_0$] (0,3) to[R,l=$R$] (3,3) to[C,l=$C$] (3,0) -- (0,0);
\draw (3,3) -- (5,3) to[L,l=$L$] (5,0) -- (3,0);
\end{circuitikz}
\caption{CircuitTikZテスト 1-1}\label{fig:circuit-1-1}
\end{figure}

\begin{boxedtheorem}{局所更新の原理}{local-update-1-1}
完全な再コンパイルではなく、構文木・ボックス・参照関係を局所的に更新できれば、TeX執筆体験はWeb DOMに近づく。
\end{boxedtheorem}

\begin{mdframed}[backgroundcolor=softblue,linecolor=mainblue,roundcorner=6pt]
\blindtext[1]
\end{mdframed}


\section{総合テストセクション 1-2}\label{sec:ch1-s2}
この節では、\TeXDOM、LuaLaTeX、数式組版、図表、相互参照、索引生成を同時に試験する。索引語としてLuaLaTeX\index{LuaLaTeX}、TikZ\index{TikZ}、tcolorbox\index{tcolorbox}、数式\index{数式}、日本語組版\index{日本語組版}を登録する。

\begin{definition}[状態ベクトル]
複素ヒルベルト空間 $\mathcal{H}$ の単位ベクトル $\ket{\psi}\in\mathcal{H}$ を純粋状態という。射影 $\rho=\ket{\psi}\bra{\psi}$ は $\rho^2=\rho$, $\Tr\rho=1$ を満たす。
\end{definition}

\begin{theorem}[スペクトル分解の簡易形]
任意のエルミート行列 $A\in\C^{n\times n}$ には、ユニタリ行列 $U$ と実対角行列 $D=\diag(\lambda_1,\dots,\lambda_n)$ が存在して
\begin{equation}\label{eq:spectral-1-2}
A=UDU^\dagger,\qquad \ee^{-\ii t A}=U\diag(\ee^{-\ii t\lambda_1},\dots,\ee^{-\ii t\lambda_n})U^\dagger
\end{equation}
と書ける。
\end{theorem}

\begin{proof}
有限次元のスペクトル定理による。相異なる固有値に属する固有空間は直交し、各固有空間で正規直交基底を選ぶことで $U$ が得られる。\cref{eq:spectral-1-2} は関数計算の定義から従う。
\end{proof}

\begin{empheq}[box=\fbox]{align}
Z(\beta)&=\Tr \ee^{-\beta H}=\sum_n \ee^{-\beta E_n},\\
F(\beta)&=-\frac{1}{\beta}\log Z(\beta),\\
\langle E\rangle&=-\frac{\partial}{\partial\beta}\log Z(\beta).
\end{empheq}

\begin{bluebox}{日本語・英語・記号混在テスト}
これは長い日本語段落である。The quick brown fox jumps over the lazy dog. 数式 $\int_{-\infty}^{\infty}\ee^{-x^2}\dd x=\sqrt{\pi}$ と、化学式 \ce{2H2 + O2 -> 2H2O} と、単位 \SI{6.62607015e-34}{\joule\second} を同時に含む。さらに、\enquote{引用符}、脚注\footnote{脚注テスト。日本語・English・$E=mc^2$ を含む。}、余白注\marginnote{Margin note 1-2}も入れる。
\end{bluebox}

\subsection{行列・テンソル・微分形式}
\begin{align}
\mat{M}_{2} &=
\begin{pmatrix}
1 & 2 & 1 \\
0 & 1 & 2+1 \\
1 & 0 & 1
\end{pmatrix},
&
\det \mat{M}_{2}&=1-1(2)(2+1),\\
\tensor{R}{^\rho_{\sigma\mu\nu}}&=\partial_\mu \Gamma^\rho_{\nu\sigma}-\partial_\nu\Gamma^\rho_{\mu\sigma}+\Gamma^\rho_{\mu\lambda}\Gamma^\lambda_{\nu\sigma}-\Gamma^\rho_{\nu\lambda}\Gamma^\lambda_{\mu\sigma},\\
\oint_{\partial\Omega}\vect{F}\cdot\vect{n}\dd S&=\iiint_{\Omega}\nabla\cdot\vect{F}\dd V.
\end{align}

\subsection{表と長いセル}
\begin{table}[H]
\centering
\caption{booktabs, makecell, multirow, colortbl の混在テスト 1-2}\label{tab:mix-1-2}
\begin{tabular}{llcr}
\toprule
\rowcolor{softblue} 分類 & パッケージ & 典型用途 & 負荷 \\
\midrule
数式 & amsmath/mathtools & \makecell{多行整列\\括弧制御} & 高 \\
図 & TikZ/PGFPlots & ベクター図形 & 高 \\
表 & longtable/tabularx & ページ分割表 & 中 \\
コード & listings/fancyvrb & ソース表示 & 中 \\
\bottomrule
\end{tabular}
\end{table}

\subsection{TikZ図}
\begin{figure}[H]
\centering
\begin{tikzpicture}[node distance=18mm,>=Latex]
\node[draw,rounded rectangle,fill=softblue,minimum width=26mm,minimum height=10mm] (src) {Source TeX};
\node[draw,rectangle,fill=softgreen,right=of src,minimum width=26mm,minimum height=10mm] (parse) {Parser};
\node[draw,diamond,aspect=2,fill=softred,right=of parse,minimum width=28mm] (dom) {TeX-DOM};
\node[draw,rounded rectangle,fill=softblue,right=of dom,minimum width=26mm,minimum height=10mm] (pdf) {PDF};
\draw[->,thick] (src) -- node[above]{tokens} (parse);
\draw[->,thick] (parse) -- node[above]{nodes} (dom);
\draw[->,thick] (dom) -- node[above]{layout} (pdf);
\draw[->,bend left=35,dashed] (pdf.north) to node[above]{SyncTeX} (src.north);
\end{tikzpicture}
\caption{TikZによる処理系模式図 1-2}\label{fig:tikz-1-2}
\end{figure}

\subsection{PGFPlots}
\begin{figure}[H]
\centering
\begin{tikzpicture}
\begin{axis}[width=.8\linewidth,height=55mm,grid=major,xlabel={$x$},ylabel={$f(x)$},legend pos=north east]
\addplot[domain=0:6.28,samples=80] {sin(deg(x))};
\addlegendentry{$\sin x$}
\addplot[domain=0:6.28,samples=80,dashed] {cos(deg(x))};
\addlegendentry{$\cos x$}
\addplot[domain=0:6.28,samples=80,dotted] {exp(-x/4)*sin(deg(3*x))};
\addlegendentry{$e^{-x/4}\sin(3x)$}
\end{axis}
\end{tikzpicture}
\caption{PGFPlotsテスト 1-2}\label{fig:pgfplots-1-2}
\end{figure}

\subsection{コードとアルゴリズム}
\begin{lstlisting}[style=pythonstyle,caption={Python風擬似コード 1-2}]
from math import exp

def partition_function(beta, energies):
    return sum(exp(-beta * e) for e in energies)

def free_energy(beta, energies):
    z = partition_function(beta, energies)
    return -log(z) / beta
\end{lstlisting}

\begin{algorithm}[H]
\caption{差分プレビュー更新アルゴリズム 1-2}
\KwIn{old DOM, new source tokens}
\KwOut{minimal changed regions}
Parse changed token interval\;
Update syntax tree nodes\;
Compute layout invalidation frontier\;
\If{math box changed}{Rebuild equation box and update cross references\;}
\Else{Reuse previous node metrics\;}
Return changed PDF rectangles\;
\end{algorithm}

\subsection{化学・回路・箱}
\begin{greenbox}{化学式と単位}
水の生成反応は \ce{2H2 + O2 -> 2H2O} であり、標準状態の議論では温度 \SI{298.15}{\kelvin}、圧力 \SI{1}{\bar} を使う。構造式の簡易表示: \chemfig{H-O-H}。
\end{greenbox}

\begin{figure}[H]
\centering
\begin{circuitikz}[scale=0.9]
\draw (0,0) to[V,l=$V_0$] (0,3) to[R,l=$R$] (3,3) to[C,l=$C$] (3,0) -- (0,0);
\draw (3,3) -- (5,3) to[L,l=$L$] (5,0) -- (3,0);
\end{circuitikz}
\caption{CircuitTikZテスト 1-2}\label{fig:circuit-1-2}
\end{figure}

\begin{boxedtheorem}{局所更新の原理}{local-update-1-2}
完全な再コンパイルではなく、構文木・ボックス・参照関係を局所的に更新できれば、TeX執筆体験はWeb DOMに近づく。
\end{boxedtheorem}

\begin{mdframed}[backgroundcolor=softblue,linecolor=mainblue,roundcorner=6pt]
\blindtext[1]
\end{mdframed}


\section{総合テストセクション 1-3}\label{sec:ch1-s3}
この節では、\TeXDOM、LuaLaTeX、数式組版、図表、相互参照、索引生成を同時に試験する。索引語としてLuaLaTeX\index{LuaLaTeX}、TikZ\index{TikZ}、tcolorbox\index{tcolorbox}、数式\index{数式}、日本語組版\index{日本語組版}を登録する。

\begin{definition}[状態ベクトル]
複素ヒルベルト空間 $\mathcal{H}$ の単位ベクトル $\ket{\psi}\in\mathcal{H}$ を純粋状態という。射影 $\rho=\ket{\psi}\bra{\psi}$ は $\rho^2=\rho$, $\Tr\rho=1$ を満たす。
\end{definition}

\begin{theorem}[スペクトル分解の簡易形]
任意のエルミート行列 $A\in\C^{n\times n}$ には、ユニタリ行列 $U$ と実対角行列 $D=\diag(\lambda_1,\dots,\lambda_n)$ が存在して
\begin{equation}\label{eq:spectral-1-3}
A=UDU^\dagger,\qquad \ee^{-\ii t A}=U\diag(\ee^{-\ii t\lambda_1},\dots,\ee^{-\ii t\lambda_n})U^\dagger
\end{equation}
と書ける。
\end{theorem}

\begin{proof}
有限次元のスペクトル定理による。相異なる固有値に属する固有空間は直交し、各固有空間で正規直交基底を選ぶことで $U$ が得られる。\cref{eq:spectral-1-3} は関数計算の定義から従う。
\end{proof}

\begin{empheq}[box=\fbox]{align}
Z(\beta)&=\Tr \ee^{-\beta H}=\sum_n \ee^{-\beta E_n},\\
F(\beta)&=-\frac{1}{\beta}\log Z(\beta),\\
\langle E\rangle&=-\frac{\partial}{\partial\beta}\log Z(\beta).
\end{empheq}

\begin{bluebox}{日本語・英語・記号混在テスト}
これは長い日本語段落である。The quick brown fox jumps over the lazy dog. 数式 $\int_{-\infty}^{\infty}\ee^{-x^2}\dd x=\sqrt{\pi}$ と、化学式 \ce{2H2 + O2 -> 2H2O} と、単位 \SI{6.62607015e-34}{\joule\second} を同時に含む。さらに、\enquote{引用符}、脚注\footnote{脚注テスト。日本語・English・$E=mc^2$ を含む。}、余白注\marginnote{Margin note 1-3}も入れる。
\end{bluebox}

\subsection{行列・テンソル・微分形式}
\begin{align}
\mat{M}_{3} &=
\begin{pmatrix}
1 & 3 & 1 \\
0 & 1 & 3+1 \\
1 & 0 & 1
\end{pmatrix},
&
\det \mat{M}_{3}&=1-1(3)(3+1),\\
\tensor{R}{^\rho_{\sigma\mu\nu}}&=\partial_\mu \Gamma^\rho_{\nu\sigma}-\partial_\nu\Gamma^\rho_{\mu\sigma}+\Gamma^\rho_{\mu\lambda}\Gamma^\lambda_{\nu\sigma}-\Gamma^\rho_{\nu\lambda}\Gamma^\lambda_{\mu\sigma},\\
\oint_{\partial\Omega}\vect{F}\cdot\vect{n}\dd S&=\iiint_{\Omega}\nabla\cdot\vect{F}\dd V.
\end{align}

\subsection{表と長いセル}
\begin{table}[H]
\centering
\caption{booktabs, makecell, multirow, colortbl の混在テスト 1-3}\label{tab:mix-1-3}
\begin{tabular}{llcr}
\toprule
\rowcolor{softblue} 分類 & パッケージ & 典型用途 & 負荷 \\
\midrule
数式 & amsmath/mathtools & \makecell{多行整列\\括弧制御} & 高 \\
図 & TikZ/PGFPlots & ベクター図形 & 高 \\
表 & longtable/tabularx & ページ分割表 & 中 \\
コード & listings/fancyvrb & ソース表示 & 中 \\
\bottomrule
\end{tabular}
\end{table}

\subsection{TikZ図}
\begin{figure}[H]
\centering
\begin{tikzpicture}[node distance=18mm,>=Latex]
\node[draw,rounded rectangle,fill=softblue,minimum width=26mm,minimum height=10mm] (src) {Source TeX};
\node[draw,rectangle,fill=softgreen,right=of src,minimum width=26mm,minimum height=10mm] (parse) {Parser};
\node[draw,diamond,aspect=2,fill=softred,right=of parse,minimum width=28mm] (dom) {TeX-DOM};
\node[draw,rounded rectangle,fill=softblue,right=of dom,minimum width=26mm,minimum height=10mm] (pdf) {PDF};
\draw[->,thick] (src) -- node[above]{tokens} (parse);
\draw[->,thick] (parse) -- node[above]{nodes} (dom);
\draw[->,thick] (dom) -- node[above]{layout} (pdf);
\draw[->,bend left=35,dashed] (pdf.north) to node[above]{SyncTeX} (src.north);
\end{tikzpicture}
\caption{TikZによる処理系模式図 1-3}\label{fig:tikz-1-3}
\end{figure}

\subsection{PGFPlots}
\begin{figure}[H]
\centering
\begin{tikzpicture}
\begin{axis}[width=.8\linewidth,height=55mm,grid=major,xlabel={$x$},ylabel={$f(x)$},legend pos=north east]
\addplot[domain=0:6.28,samples=80] {sin(deg(x))};
\addlegendentry{$\sin x$}
\addplot[domain=0:6.28,samples=80,dashed] {cos(deg(x))};
\addlegendentry{$\cos x$}
\addplot[domain=0:6.28,samples=80,dotted] {exp(-x/4)*sin(deg(3*x))};
\addlegendentry{$e^{-x/4}\sin(3x)$}
\end{axis}
\end{tikzpicture}
\caption{PGFPlotsテスト 1-3}\label{fig:pgfplots-1-3}
\end{figure}

\subsection{コードとアルゴリズム}
\begin{lstlisting}[style=pythonstyle,caption={Python風擬似コード 1-3}]
from math import exp

def partition_function(beta, energies):
    return sum(exp(-beta * e) for e in energies)

def free_energy(beta, energies):
    z = partition_function(beta, energies)
    return -log(z) / beta
\end{lstlisting}

\begin{algorithm}[H]
\caption{差分プレビュー更新アルゴリズム 1-3}
\KwIn{old DOM, new source tokens}
\KwOut{minimal changed regions}
Parse changed token interval\;
Update syntax tree nodes\;
Compute layout invalidation frontier\;
\If{math box changed}{Rebuild equation box and update cross references\;}
\Else{Reuse previous node metrics\;}
Return changed PDF rectangles\;
\end{algorithm}

\subsection{化学・回路・箱}
\begin{greenbox}{化学式と単位}
水の生成反応は \ce{2H2 + O2 -> 2H2O} であり、標準状態の議論では温度 \SI{298.15}{\kelvin}、圧力 \SI{1}{\bar} を使う。構造式の簡易表示: \chemfig{H-O-H}。
\end{greenbox}

\begin{figure}[H]
\centering
\begin{circuitikz}[scale=0.9]
\draw (0,0) to[V,l=$V_0$] (0,3) to[R,l=$R$] (3,3) to[C,l=$C$] (3,0) -- (0,0);
\draw (3,3) -- (5,3) to[L,l=$L$] (5,0) -- (3,0);
\end{circuitikz}
\caption{CircuitTikZテスト 1-3}\label{fig:circuit-1-3}
\end{figure}

\begin{boxedtheorem}{局所更新の原理}{local-update-1-3}
完全な再コンパイルではなく、構文木・ボックス・参照関係を局所的に更新できれば、TeX執筆体験はWeb DOMに近づく。
\end{boxedtheorem}

\begin{mdframed}[backgroundcolor=softblue,linecolor=mainblue,roundcorner=6pt]
\blindtext[1]
\end{mdframed}

\chapter{巨大総合章 2: パッケージ横断テスト}\label{chap:2}


\begin{redbox}{章の目的}
この章は、同じパターンを少しずつ変えながら反復し、アウトライン、ラベル、図表番号、ページ分割、脚注、索引、長い数式、コード、図形描画を大量に発生させるための章である。
\end{redbox}


\section{総合テストセクション 2-1}\label{sec:ch2-s1}
この節では、\TeXDOM、LuaLaTeX、数式組版、図表、相互参照、索引生成を同時に試験する。索引語としてLuaLaTeX\index{LuaLaTeX}、TikZ\index{TikZ}、tcolorbox\index{tcolorbox}、数式\index{数式}、日本語組版\index{日本語組版}を登録する。

\begin{definition}[状態ベクトル]
複素ヒルベルト空間 $\mathcal{H}$ の単位ベクトル $\ket{\psi}\in\mathcal{H}$ を純粋状態という。射影 $\rho=\ket{\psi}\bra{\psi}$ は $\rho^2=\rho$, $\Tr\rho=1$ を満たす。
\end{definition}

\begin{theorem}[スペクトル分解の簡易形]
任意のエルミート行列 $A\in\C^{n\times n}$ には、ユニタリ行列 $U$ と実対角行列 $D=\diag(\lambda_1,\dots,\lambda_n)$ が存在して
\begin{equation}\label{eq:spectral-2-1}
A=UDU^\dagger,\qquad \ee^{-\ii t A}=U\diag(\ee^{-\ii t\lambda_1},\dots,\ee^{-\ii t\lambda_n})U^\dagger
\end{equation}
と書ける。
\end{theorem}

\begin{proof}
有限次元のスペクトル定理による。相異なる固有値に属する固有空間は直交し、各固有空間で正規直交基底を選ぶことで $U$ が得られる。\cref{eq:spectral-2-1} は関数計算の定義から従う。
\end{proof}

\begin{empheq}[box=\fbox]{align}
Z(\beta)&=\Tr \ee^{-\beta H}=\sum_n \ee^{-\beta E_n},\\
F(\beta)&=-\frac{1}{\beta}\log Z(\beta),\\
\langle E\rangle&=-\frac{\partial}{\partial\beta}\log Z(\beta).
\end{empheq}

\begin{bluebox}{日本語・英語・記号混在テスト}
これは長い日本語段落である。The quick brown fox jumps over the lazy dog. 数式 $\int_{-\infty}^{\infty}\ee^{-x^2}\dd x=\sqrt{\pi}$ と、化学式 \ce{2H2 + O2 -> 2H2O} と、単位 \SI{6.62607015e-34}{\joule\second} を同時に含む。さらに、\enquote{引用符}、脚注\footnote{脚注テスト。日本語・English・$E=mc^2$ を含む。}、余白注\marginnote{Margin note 2-1}も入れる。
\end{bluebox}

\subsection{行列・テンソル・微分形式}
\begin{align}
\mat{M}_{1} &=
\begin{pmatrix}
1 & 1 & 2 \\
0 & 1 & 1+2 \\
2 & 0 & 1
\end{pmatrix},
&
\det \mat{M}_{1}&=1-2(1)(1+2),\\
\tensor{R}{^\rho_{\sigma\mu\nu}}&=\partial_\mu \Gamma^\rho_{\nu\sigma}-\partial_\nu\Gamma^\rho_{\mu\sigma}+\Gamma^\rho_{\mu\lambda}\Gamma^\lambda_{\nu\sigma}-\Gamma^\rho_{\nu\lambda}\Gamma^\lambda_{\mu\sigma},\\
\oint_{\partial\Omega}\vect{F}\cdot\vect{n}\dd S&=\iiint_{\Omega}\nabla\cdot\vect{F}\dd V.
\end{align}

\subsection{表と長いセル}
\begin{table}[H]
\centering
\caption{booktabs, makecell, multirow, colortbl の混在テスト 2-1}\label{tab:mix-2-1}
\begin{tabular}{llcr}
\toprule
\rowcolor{softblue} 分類 & パッケージ & 典型用途 & 負荷 \\
\midrule
数式 & amsmath/mathtools & \makecell{多行整列\\括弧制御} & 高 \\
図 & TikZ/PGFPlots & ベクター図形 & 高 \\
表 & longtable/tabularx & ページ分割表 & 中 \\
コード & listings/fancyvrb & ソース表示 & 中 \\
\bottomrule
\end{tabular}
\end{table}

\subsection{TikZ図}
\begin{figure}[H]
\centering
\begin{tikzpicture}[node distance=18mm,>=Latex]
\node[draw,rounded rectangle,fill=softblue,minimum width=26mm,minimum height=10mm] (src) {Source TeX};
\node[draw,rectangle,fill=softgreen,right=of src,minimum width=26mm,minimum height=10mm] (parse) {Parser};
\node[draw,diamond,aspect=2,fill=softred,right=of parse,minimum width=28mm] (dom) {TeX-DOM};
\node[draw,rounded rectangle,fill=softblue,right=of dom,minimum width=26mm,minimum height=10mm] (pdf) {PDF};
\draw[->,thick] (src) -- node[above]{tokens} (parse);
\draw[->,thick] (parse) -- node[above]{nodes} (dom);
\draw[->,thick] (dom) -- node[above]{layout} (pdf);
\draw[->,bend left=35,dashed] (pdf.north) to node[above]{SyncTeX} (src.north);
\end{tikzpicture}
\caption{TikZによる処理系模式図 2-1}\label{fig:tikz-2-1}
\end{figure}

\subsection{PGFPlots}
\begin{figure}[H]
\centering
\begin{tikzpicture}
\begin{axis}[width=.8\linewidth,height=55mm,grid=major,xlabel={$x$},ylabel={$f(x)$},legend pos=north east]
\addplot[domain=0:6.28,samples=80] {sin(deg(x))};
\addlegendentry{$\sin x$}
\addplot[domain=0:6.28,samples=80,dashed] {cos(deg(x))};
\addlegendentry{$\cos x$}
\addplot[domain=0:6.28,samples=80,dotted] {exp(-x/4)*sin(deg(3*x))};
\addlegendentry{$e^{-x/4}\sin(3x)$}
\end{axis}
\end{tikzpicture}
\caption{PGFPlotsテスト 2-1}\label{fig:pgfplots-2-1}
\end{figure}

\subsection{コードとアルゴリズム}
\begin{lstlisting}[style=pythonstyle,caption={Python風擬似コード 2-1}]
from math import exp

def partition_function(beta, energies):
    return sum(exp(-beta * e) for e in energies)

def free_energy(beta, energies):
    z = partition_function(beta, energies)
    return -log(z) / beta
\end{lstlisting}

\begin{algorithm}[H]
\caption{差分プレビュー更新アルゴリズム 2-1}
\KwIn{old DOM, new source tokens}
\KwOut{minimal changed regions}
Parse changed token interval\;
Update syntax tree nodes\;
Compute layout invalidation frontier\;
\If{math box changed}{Rebuild equation box and update cross references\;}
\Else{Reuse previous node metrics\;}
Return changed PDF rectangles\;
\end{algorithm}

\subsection{化学・回路・箱}
\begin{greenbox}{化学式と単位}
水の生成反応は \ce{2H2 + O2 -> 2H2O} であり、標準状態の議論では温度 \SI{298.15}{\kelvin}、圧力 \SI{1}{\bar} を使う。構造式の簡易表示: \chemfig{H-O-H}。
\end{greenbox}

\begin{figure}[H]
\centering
\begin{circuitikz}[scale=0.9]
\draw (0,0) to[V,l=$V_0$] (0,3) to[R,l=$R$] (3,3) to[C,l=$C$] (3,0) -- (0,0);
\draw (3,3) -- (5,3) to[L,l=$L$] (5,0) -- (3,0);
\end{circuitikz}
\caption{CircuitTikZテスト 2-1}\label{fig:circuit-2-1}
\end{figure}

\begin{boxedtheorem}{局所更新の原理}{local-update-2-1}
完全な再コンパイルではなく、構文木・ボックス・参照関係を局所的に更新できれば、TeX執筆体験はWeb DOMに近づく。
\end{boxedtheorem}

\begin{mdframed}[backgroundcolor=softblue,linecolor=mainblue,roundcorner=6pt]
\blindtext[1]
\end{mdframed}


\section{総合テストセクション 2-2}\label{sec:ch2-s2}
この節では、\TeXDOM、LuaLaTeX、数式組版、図表、相互参照、索引生成を同時に試験する。索引語としてLuaLaTeX\index{LuaLaTeX}、TikZ\index{TikZ}、tcolorbox\index{tcolorbox}、数式\index{数式}、日本語組版\index{日本語組版}を登録する。

\begin{definition}[状態ベクトル]
複素ヒルベルト空間 $\mathcal{H}$ の単位ベクトル $\ket{\psi}\in\mathcal{H}$ を純粋状態という。射影 $\rho=\ket{\psi}\bra{\psi}$ は $\rho^2=\rho$, $\Tr\rho=1$ を満たす。
\end{definition}

\begin{theorem}[スペクトル分解の簡易形]
任意のエルミート行列 $A\in\C^{n\times n}$ には、ユニタリ行列 $U$ と実対角行列 $D=\diag(\lambda_1,\dots,\lambda_n)$ が存在して
\begin{equation}\label{eq:spectral-2-2}
A=UDU^\dagger,\qquad \ee^{-\ii t A}=U\diag(\ee^{-\ii t\lambda_1},\dots,\ee^{-\ii t\lambda_n})U^\dagger
\end{equation}
と書ける。
\end{theorem}

\begin{proof}
有限次元のスペクトル定理による。相異なる固有値に属する固有空間は直交し、各固有空間で正規直交基底を選ぶことで $U$ が得られる。\cref{eq:spectral-2-2} は関数計算の定義から従う。
\end{proof}

\begin{empheq}[box=\fbox]{align}
Z(\beta)&=\Tr \ee^{-\beta H}=\sum_n \ee^{-\beta E_n},\\
F(\beta)&=-\frac{1}{\beta}\log Z(\beta),\\
\langle E\rangle&=-\frac{\partial}{\partial\beta}\log Z(\beta).
\end{empheq}

\begin{bluebox}{日本語・英語・記号混在テスト}
これは長い日本語段落である。The quick brown fox jumps over the lazy dog. 数式 $\int_{-\infty}^{\infty}\ee^{-x^2}\dd x=\sqrt{\pi}$ と、化学式 \ce{2H2 + O2 -> 2H2O} と、単位 \SI{6.62607015e-34}{\joule\second} を同時に含む。さらに、\enquote{引用符}、脚注\footnote{脚注テスト。日本語・English・$E=mc^2$ を含む。}、余白注\marginnote{Margin note 2-2}も入れる。
\end{bluebox}

\subsection{行列・テンソル・微分形式}
\begin{align}
\mat{M}_{2} &=
\begin{pmatrix}
1 & 2 & 2 \\
0 & 1 & 2+2 \\
2 & 0 & 1
\end{pmatrix},
&
\det \mat{M}_{2}&=1-2(2)(2+2),\\
\tensor{R}{^\rho_{\sigma\mu\nu}}&=\partial_\mu \Gamma^\rho_{\nu\sigma}-\partial_\nu\Gamma^\rho_{\mu\sigma}+\Gamma^\rho_{\mu\lambda}\Gamma^\lambda_{\nu\sigma}-\Gamma^\rho_{\nu\lambda}\Gamma^\lambda_{\mu\sigma},\\
\oint_{\partial\Omega}\vect{F}\cdot\vect{n}\dd S&=\iiint_{\Omega}\nabla\cdot\vect{F}\dd V.
\end{align}

\subsection{表と長いセル}
\begin{table}[H]
\centering
\caption{booktabs, makecell, multirow, colortbl の混在テスト 2-2}\label{tab:mix-2-2}
\begin{tabular}{llcr}
\toprule
\rowcolor{softblue} 分類 & パッケージ & 典型用途 & 負荷 \\
\midrule
数式 & amsmath/mathtools & \makecell{多行整列\\括弧制御} & 高 \\
図 & TikZ/PGFPlots & ベクター図形 & 高 \\
表 & longtable/tabularx & ページ分割表 & 中 \\
コード & listings/fancyvrb & ソース表示 & 中 \\
\bottomrule
\end{tabular}
\end{table}

\subsection{TikZ図}
\begin{figure}[H]
\centering
\begin{tikzpicture}[node distance=18mm,>=Latex]
\node[draw,rounded rectangle,fill=softblue,minimum width=26mm,minimum height=10mm] (src) {Source TeX};
\node[draw,rectangle,fill=softgreen,right=of src,minimum width=26mm,minimum height=10mm] (parse) {Parser};
\node[draw,diamond,aspect=2,fill=softred,right=of parse,minimum width=28mm] (dom) {TeX-DOM};
\node[draw,rounded rectangle,fill=softblue,right=of dom,minimum width=26mm,minimum height=10mm] (pdf) {PDF};
\draw[->,thick] (src) -- node[above]{tokens} (parse);
\draw[->,thick] (parse) -- node[above]{nodes} (dom);
\draw[->,thick] (dom) -- node[above]{layout} (pdf);
\draw[->,bend left=35,dashed] (pdf.north) to node[above]{SyncTeX} (src.north);
\end{tikzpicture}
\caption{TikZによる処理系模式図 2-2}\label{fig:tikz-2-2}
\end{figure}

\subsection{PGFPlots}
\begin{figure}[H]
\centering
\begin{tikzpicture}
\begin{axis}[width=.8\linewidth,height=55mm,grid=major,xlabel={$x$},ylabel={$f(x)$},legend pos=north east]
\addplot[domain=0:6.28,samples=80] {sin(deg(x))};
\addlegendentry{$\sin x$}
\addplot[domain=0:6.28,samples=80,dashed] {cos(deg(x))};
\addlegendentry{$\cos x$}
\addplot[domain=0:6.28,samples=80,dotted] {exp(-x/4)*sin(deg(3*x))};
\addlegendentry{$e^{-x/4}\sin(3x)$}
\end{axis}
\end{tikzpicture}
\caption{PGFPlotsテスト 2-2}\label{fig:pgfplots-2-2}
\end{figure}

\subsection{コードとアルゴリズム}
\begin{lstlisting}[style=pythonstyle,caption={Python風擬似コード 2-2}]
from math import exp

def partition_function(beta, energies):
    return sum(exp(-beta * e) for e in energies)

def free_energy(beta, energies):
    z = partition_function(beta, energies)
    return -log(z) / beta
\end{lstlisting}

\begin{algorithm}[H]
\caption{差分プレビュー更新アルゴリズム 2-2}
\KwIn{old DOM, new source tokens}
\KwOut{minimal changed regions}
Parse changed token interval\;
Update syntax tree nodes\;
Compute layout invalidation frontier\;
\If{math box changed}{Rebuild equation box and update cross references\;}
\Else{Reuse previous node metrics\;}
Return changed PDF rectangles\;
\end{algorithm}

\subsection{化学・回路・箱}
\begin{greenbox}{化学式と単位}
水の生成反応は \ce{2H2 + O2 -> 2H2O} であり、標準状態の議論では温度 \SI{298.15}{\kelvin}、圧力 \SI{1}{\bar} を使う。構造式の簡易表示: \chemfig{H-O-H}。
\end{greenbox}

\begin{figure}[H]
\centering
\begin{circuitikz}[scale=0.9]
\draw (0,0) to[V,l=$V_0$] (0,3) to[R,l=$R$] (3,3) to[C,l=$C$] (3,0) -- (0,0);
\draw (3,3) -- (5,3) to[L,l=$L$] (5,0) -- (3,0);
\end{circuitikz}
\caption{CircuitTikZテスト 2-2}\label{fig:circuit-2-2}
\end{figure}

\begin{boxedtheorem}{局所更新の原理}{local-update-2-2}
完全な再コンパイルではなく、構文木・ボックス・参照関係を局所的に更新できれば、TeX執筆体験はWeb DOMに近づく。
\end{boxedtheorem}

\begin{mdframed}[backgroundcolor=softblue,linecolor=mainblue,roundcorner=6pt]
\blindtext[1]
\end{mdframed}


\section{総合テストセクション 2-3}\label{sec:ch2-s3}
この節では、\TeXDOM、LuaLaTeX、数式組版、図表、相互参照、索引生成を同時に試験する。索引語としてLuaLaTeX\index{LuaLaTeX}、TikZ\index{TikZ}、tcolorbox\index{tcolorbox}、数式\index{数式}、日本語組版\index{日本語組版}を登録する。

\begin{definition}[状態ベクトル]
複素ヒルベルト空間 $\mathcal{H}$ の単位ベクトル $\ket{\psi}\in\mathcal{H}$ を純粋状態という。射影 $\rho=\ket{\psi}\bra{\psi}$ は $\rho^2=\rho$, $\Tr\rho=1$ を満たす。
\end{definition}

\begin{theorem}[スペクトル分解の簡易形]
任意のエルミート行列 $A\in\C^{n\times n}$ には、ユニタリ行列 $U$ と実対角行列 $D=\diag(\lambda_1,\dots,\lambda_n)$ が存在して
\begin{equation}\label{eq:spectral-2-3}
A=UDU^\dagger,\qquad \ee^{-\ii t A}=U\diag(\ee^{-\ii t\lambda_1},\dots,\ee^{-\ii t\lambda_n})U^\dagger
\end{equation}
と書ける。
\end{theorem}

\begin{proof}
有限次元のスペクトル定理による。相異なる固有値に属する固有空間は直交し、各固有空間で正規直交基底を選ぶことで $U$ が得られる。\cref{eq:spectral-2-3} は関数計算の定義から従う。
\end{proof}

\begin{empheq}[box=\fbox]{align}
Z(\beta)&=\Tr \ee^{-\beta H}=\sum_n \ee^{-\beta E_n},\\
F(\beta)&=-\frac{1}{\beta}\log Z(\beta),\\
\langle E\rangle&=-\frac{\partial}{\partial\beta}\log Z(\beta).
\end{empheq}

\begin{bluebox}{日本語・英語・記号混在テスト}
これは長い日本語段落である。The quick brown fox jumps over the lazy dog. 数式 $\int_{-\infty}^{\infty}\ee^{-x^2}\dd x=\sqrt{\pi}$ と、化学式 \ce{2H2 + O2 -> 2H2O} と、単位 \SI{6.62607015e-34}{\joule\second} を同時に含む。さらに、\enquote{引用符}、脚注\footnote{脚注テスト。日本語・English・$E=mc^2$ を含む。}、余白注\marginnote{Margin note 2-3}も入れる。
\end{bluebox}

\subsection{行列・テンソル・微分形式}
\begin{align}
\mat{M}_{3} &=
\begin{pmatrix}
1 & 3 & 2 \\
0 & 1 & 3+2 \\
2 & 0 & 1
\end{pmatrix},
&
\det \mat{M}_{3}&=1-2(3)(3+2),\\
\tensor{R}{^\rho_{\sigma\mu\nu}}&=\partial_\mu \Gamma^\rho_{\nu\sigma}-\partial_\nu\Gamma^\rho_{\mu\sigma}+\Gamma^\rho_{\mu\lambda}\Gamma^\lambda_{\nu\sigma}-\Gamma^\rho_{\nu\lambda}\Gamma^\lambda_{\mu\sigma},\\
\oint_{\partial\Omega}\vect{F}\cdot\vect{n}\dd S&=\iiint_{\Omega}\nabla\cdot\vect{F}\dd V.
\end{align}

\subsection{表と長いセル}
\begin{table}[H]
\centering
\caption{booktabs, makecell, multirow, colortbl の混在テスト 2-3}\label{tab:mix-2-3}
\begin{tabular}{llcr}
\toprule
\rowcolor{softblue} 分類 & パッケージ & 典型用途 & 負荷 \\
\midrule
数式 & amsmath/mathtools & \makecell{多行整列\\括弧制御} & 高 \\
図 & TikZ/PGFPlots & ベクター図形 & 高 \\
表 & longtable/tabularx & ページ分割表 & 中 \\
コード & listings/fancyvrb & ソース表示 & 中 \\
\bottomrule
\end{tabular}
\end{table}

\subsection{TikZ図}
\begin{figure}[H]
\centering
\begin{tikzpicture}[node distance=18mm,>=Latex]
\node[draw,rounded rectangle,fill=softblue,minimum width=26mm,minimum height=10mm] (src) {Source TeX};
\node[draw,rectangle,fill=softgreen,right=of src,minimum width=26mm,minimum height=10mm] (parse) {Parser};
\node[draw,diamond,aspect=2,fill=softred,right=of parse,minimum width=28mm] (dom) {TeX-DOM};
\node[draw,rounded rectangle,fill=softblue,right=of dom,minimum width=26mm,minimum height=10mm] (pdf) {PDF};
\draw[->,thick] (src) -- node[above]{tokens} (parse);
\draw[->,thick] (parse) -- node[above]{nodes} (dom);
\draw[->,thick] (dom) -- node[above]{layout} (pdf);
\draw[->,bend left=35,dashed] (pdf.north) to node[above]{SyncTeX} (src.north);
\end{tikzpicture}
\caption{TikZによる処理系模式図 2-3}\label{fig:tikz-2-3}
\end{figure}

\subsection{PGFPlots}
\begin{figure}[H]
\centering
\begin{tikzpicture}
\begin{axis}[width=.8\linewidth,height=55mm,grid=major,xlabel={$x$},ylabel={$f(x)$},legend pos=north east]
\addplot[domain=0:6.28,samples=80] {sin(deg(x))};
\addlegendentry{$\sin x$}
\addplot[domain=0:6.28,samples=80,dashed] {cos(deg(x))};
\addlegendentry{$\cos x$}
\addplot[domain=0:6.28,samples=80,dotted] {exp(-x/4)*sin(deg(3*x))};
\addlegendentry{$e^{-x/4}\sin(3x)$}
\end{axis}
\end{tikzpicture}
\caption{PGFPlotsテスト 2-3}\label{fig:pgfplots-2-3}
\end{figure}

\subsection{コードとアルゴリズム}
\begin{lstlisting}[style=pythonstyle,caption={Python風擬似コード 2-3}]
from math import exp

def partition_function(beta, energies):
    return sum(exp(-beta * e) for e in energies)

def free_energy(beta, energies):
    z = partition_function(beta, energies)
    return -log(z) / beta
\end{lstlisting}

\begin{algorithm}[H]
\caption{差分プレビュー更新アルゴリズム 2-3}
\KwIn{old DOM, new source tokens}
\KwOut{minimal changed regions}
Parse changed token interval\;
Update syntax tree nodes\;
Compute layout invalidation frontier\;
\If{math box changed}{Rebuild equation box and update cross references\;}
\Else{Reuse previous node metrics\;}
Return changed PDF rectangles\;
\end{algorithm}

\subsection{化学・回路・箱}
\begin{greenbox}{化学式と単位}
水の生成反応は \ce{2H2 + O2 -> 2H2O} であり、標準状態の議論では温度 \SI{298.15}{\kelvin}、圧力 \SI{1}{\bar} を使う。構造式の簡易表示: \chemfig{H-O-H}。
\end{greenbox}

\begin{figure}[H]
\centering
\begin{circuitikz}[scale=0.9]
\draw (0,0) to[V,l=$V_0$] (0,3) to[R,l=$R$] (3,3) to[C,l=$C$] (3,0) -- (0,0);
\draw (3,3) -- (5,3) to[L,l=$L$] (5,0) -- (3,0);
\end{circuitikz}
\caption{CircuitTikZテスト 2-3}\label{fig:circuit-2-3}
\end{figure}

\begin{boxedtheorem}{局所更新の原理}{local-update-2-3}
完全な再コンパイルではなく、構文木・ボックス・参照関係を局所的に更新できれば、TeX執筆体験はWeb DOMに近づく。
\end{boxedtheorem}

\begin{mdframed}[backgroundcolor=softblue,linecolor=mainblue,roundcorner=6pt]
\blindtext[1]
\end{mdframed}

\chapter{巨大総合章 3: パッケージ横断テスト}\label{chap:3}


\begin{redbox}{章の目的}
この章は、同じパターンを少しずつ変えながら反復し、アウトライン、ラベル、図表番号、ページ分割、脚注、索引、長い数式、コード、図形描画を大量に発生させるための章である。
\end{redbox}


\section{総合テストセクション 3-1}\label{sec:ch3-s1}
この節では、\TeXDOM、LuaLaTeX、数式組版、図表、相互参照、索引生成を同時に試験する。索引語としてLuaLaTeX\index{LuaLaTeX}、TikZ\index{TikZ}、tcolorbox\index{tcolorbox}、数式\index{数式}、日本語組版\index{日本語組版}を登録する。

\begin{definition}[状態ベクトル]
複素ヒルベルト空間 $\mathcal{H}$ の単位ベクトル $\ket{\psi}\in\mathcal{H}$ を純粋状態という。射影 $\rho=\ket{\psi}\bra{\psi}$ は $\rho^2=\rho$, $\Tr\rho=1$ を満たす。
\end{definition}

\begin{theorem}[スペクトル分解の簡易形]
任意のエルミート行列 $A\in\C^{n\times n}$ には、ユニタリ行列 $U$ と実対角行列 $D=\diag(\lambda_1,\dots,\lambda_n)$ が存在して
\begin{equation}\label{eq:spectral-3-1}
A=UDU^\dagger,\qquad \ee^{-\ii t A}=U\diag(\ee^{-\ii t\lambda_1},\dots,\ee^{-\ii t\lambda_n})U^\dagger
\end{equation}
と書ける。
\end{theorem}

\begin{proof}
有限次元のスペクトル定理による。相異なる固有値に属する固有空間は直交し、各固有空間で正規直交基底を選ぶことで $U$ が得られる。\cref{eq:spectral-3-1} は関数計算の定義から従う。
\end{proof}

\begin{empheq}[box=\fbox]{align}
Z(\beta)&=\Tr \ee^{-\beta H}=\sum_n \ee^{-\beta E_n},\\
F(\beta)&=-\frac{1}{\beta}\log Z(\beta),\\
\langle E\rangle&=-\frac{\partial}{\partial\beta}\log Z(\beta).
\end{empheq}

\begin{bluebox}{日本語・英語・記号混在テスト}
これは長い日本語段落である。The quick brown fox jumps over the lazy dog. 数式 $\int_{-\infty}^{\infty}\ee^{-x^2}\dd x=\sqrt{\pi}$ と、化学式 \ce{2H2 + O2 -> 2H2O} と、単位 \SI{6.62607015e-34}{\joule\second} を同時に含む。さらに、\enquote{引用符}、脚注\footnote{脚注テスト。日本語・English・$E=mc^2$ を含む。}、余白注\marginnote{Margin note 3-1}も入れる。
\end{bluebox}

\subsection{行列・テンソル・微分形式}
\begin{align}
\mat{M}_{1} &=
\begin{pmatrix}
1 & 1 & 3 \\
0 & 1 & 1+3 \\
3 & 0 & 1
\end{pmatrix},
&
\det \mat{M}_{1}&=1-3(1)(1+3),\\
\tensor{R}{^\rho_{\sigma\mu\nu}}&=\partial_\mu \Gamma^\rho_{\nu\sigma}-\partial_\nu\Gamma^\rho_{\mu\sigma}+\Gamma^\rho_{\mu\lambda}\Gamma^\lambda_{\nu\sigma}-\Gamma^\rho_{\nu\lambda}\Gamma^\lambda_{\mu\sigma},\\
\oint_{\partial\Omega}\vect{F}\cdot\vect{n}\dd S&=\iiint_{\Omega}\nabla\cdot\vect{F}\dd V.
\end{align}

\subsection{表と長いセル}
\begin{table}[H]
\centering
\caption{booktabs, makecell, multirow, colortbl の混在テスト 3-1}\label{tab:mix-3-1}
\begin{tabular}{llcr}
\toprule
\rowcolor{softblue} 分類 & パッケージ & 典型用途 & 負荷 \\
\midrule
数式 & amsmath/mathtools & \makecell{多行整列\\括弧制御} & 高 \\
図 & TikZ/PGFPlots & ベクター図形 & 高 \\
表 & longtable/tabularx & ページ分割表 & 中 \\
コード & listings/fancyvrb & ソース表示 & 中 \\
\bottomrule
\end{tabular}
\end{table}

\subsection{TikZ図}
\begin{figure}[H]
\centering
\begin{tikzpicture}[node distance=18mm,>=Latex]
\node[draw,rounded rectangle,fill=softblue,minimum width=26mm,minimum height=10mm] (src) {Source TeX};
\node[draw,rectangle,fill=softgreen,right=of src,minimum width=26mm,minimum height=10mm] (parse) {Parser};
\node[draw,diamond,aspect=2,fill=softred,right=of parse,minimum width=28mm] (dom) {TeX-DOM};
\node[draw,rounded rectangle,fill=softblue,right=of dom,minimum width=26mm,minimum height=10mm] (pdf) {PDF};
\draw[->,thick] (src) -- node[above]{tokens} (parse);
\draw[->,thick] (parse) -- node[above]{nodes} (dom);
\draw[->,thick] (dom) -- node[above]{layout} (pdf);
\draw[->,bend left=35,dashed] (pdf.north) to node[above]{SyncTeX} (src.north);
\end{tikzpicture}
\caption{TikZによる処理系模式図 3-1}\label{fig:tikz-3-1}
\end{figure}

\subsection{PGFPlots}
\begin{figure}[H]
\centering
\begin{tikzpicture}
\begin{axis}[width=.8\linewidth,height=55mm,grid=major,xlabel={$x$},ylabel={$f(x)$},legend pos=north east]
\addplot[domain=0:6.28,samples=80] {sin(deg(x))};
\addlegendentry{$\sin x$}
\addplot[domain=0:6.28,samples=80,dashed] {cos(deg(x))};
\addlegendentry{$\cos x$}
\addplot[domain=0:6.28,samples=80,dotted] {exp(-x/4)*sin(deg(3*x))};
\addlegendentry{$e^{-x/4}\sin(3x)$}
\end{axis}
\end{tikzpicture}
\caption{PGFPlotsテスト 3-1}\label{fig:pgfplots-3-1}
\end{figure}

\subsection{コードとアルゴリズム}
\begin{lstlisting}[style=pythonstyle,caption={Python風擬似コード 3-1}]
from math import exp

def partition_function(beta, energies):
    return sum(exp(-beta * e) for e in energies)

def free_energy(beta, energies):
    z = partition_function(beta, energies)
    return -log(z) / beta
\end{lstlisting}

\begin{algorithm}[H]
\caption{差分プレビュー更新アルゴリズム 3-1}
\KwIn{old DOM, new source tokens}
\KwOut{minimal changed regions}
Parse changed token interval\;
Update syntax tree nodes\;
Compute layout invalidation frontier\;
\If{math box changed}{Rebuild equation box and update cross references\;}
\Else{Reuse previous node metrics\;}
Return changed PDF rectangles\;
\end{algorithm}

\subsection{化学・回路・箱}
\begin{greenbox}{化学式と単位}
水の生成反応は \ce{2H2 + O2 -> 2H2O} であり、標準状態の議論では温度 \SI{298.15}{\kelvin}、圧力 \SI{1}{\bar} を使う。構造式の簡易表示: \chemfig{H-O-H}。
\end{greenbox}

\begin{figure}[H]
\centering
\begin{circuitikz}[scale=0.9]
\draw (0,0) to[V,l=$V_0$] (0,3) to[R,l=$R$] (3,3) to[C,l=$C$] (3,0) -- (0,0);
\draw (3,3) -- (5,3) to[L,l=$L$] (5,0) -- (3,0);
\end{circuitikz}
\caption{CircuitTikZテスト 3-1}\label{fig:circuit-3-1}
\end{figure}

\begin{boxedtheorem}{局所更新の原理}{local-update-3-1}
完全な再コンパイルではなく、構文木・ボックス・参照関係を局所的に更新できれば、TeX執筆体験はWeb DOMに近づく。
\end{boxedtheorem}

\begin{mdframed}[backgroundcolor=softblue,linecolor=mainblue,roundcorner=6pt]
\blindtext[1]
\end{mdframed}


\section{総合テストセクション 3-2}\label{sec:ch3-s2}
この節では、\TeXDOM、LuaLaTeX、数式組版、図表、相互参照、索引生成を同時に試験する。索引語としてLuaLaTeX\index{LuaLaTeX}、TikZ\index{TikZ}、tcolorbox\index{tcolorbox}、数式\index{数式}、日本語組版\index{日本語組版}を登録する。

\begin{definition}[状態ベクトル]
複素ヒルベルト空間 $\mathcal{H}$ の単位ベクトル $\ket{\psi}\in\mathcal{H}$ を純粋状態という。射影 $\rho=\ket{\psi}\bra{\psi}$ は $\rho^2=\rho$, $\Tr\rho=1$ を満たす。
\end{definition}

\begin{theorem}[スペクトル分解の簡易形]
任意のエルミート行列 $A\in\C^{n\times n}$ には、ユニタリ行列 $U$ と実対角行列 $D=\diag(\lambda_1,\dots,\lambda_n)$ が存在して
\begin{equation}\label{eq:spectral-3-2}
A=UDU^\dagger,\qquad \ee^{-\ii t A}=U\diag(\ee^{-\ii t\lambda_1},\dots,\ee^{-\ii t\lambda_n})U^\dagger
\end{equation}
と書ける。
\end{theorem}

\begin{proof}
有限次元のスペクトル定理による。相異なる固有値に属する固有空間は直交し、各固有空間で正規直交基底を選ぶことで $U$ が得られる。\cref{eq:spectral-3-2} は関数計算の定義から従う。
\end{proof}

\begin{empheq}[box=\fbox]{align}
Z(\beta)&=\Tr \ee^{-\beta H}=\sum_n \ee^{-\beta E_n},\\
F(\beta)&=-\frac{1}{\beta}\log Z(\beta),\\
\langle E\rangle&=-\frac{\partial}{\partial\beta}\log Z(\beta).
\end{empheq}

\begin{bluebox}{日本語・英語・記号混在テスト}
これは長い日本語段落である。The quick brown fox jumps over the lazy dog. 数式 $\int_{-\infty}^{\infty}\ee^{-x^2}\dd x=\sqrt{\pi}$ と、化学式 \ce{2H2 + O2 -> 2H2O} と、単位 \SI{6.62607015e-34}{\joule\second} を同時に含む。さらに、\enquote{引用符}、脚注\footnote{脚注テスト。日本語・English・$E=mc^2$ を含む。}、余白注\marginnote{Margin note 3-2}も入れる。
\end{bluebox}

\subsection{行列・テンソル・微分形式}
\begin{align}
\mat{M}_{2} &=
\begin{pmatrix}
1 & 2 & 3 \\
0 & 1 & 2+3 \\
3 & 0 & 1
\end{pmatrix},
&
\det \mat{M}_{2}&=1-3(2)(2+3),\\
\tensor{R}{^\rho_{\sigma\mu\nu}}&=\partial_\mu \Gamma^\rho_{\nu\sigma}-\partial_\nu\Gamma^\rho_{\mu\sigma}+\Gamma^\rho_{\mu\lambda}\Gamma^\lambda_{\nu\sigma}-\Gamma^\rho_{\nu\lambda}\Gamma^\lambda_{\mu\sigma},\\
\oint_{\partial\Omega}\vect{F}\cdot\vect{n}\dd S&=\iiint_{\Omega}\nabla\cdot\vect{F}\dd V.
\end{align}

\subsection{表と長いセル}
\begin{table}[H]
\centering
\caption{booktabs, makecell, multirow, colortbl の混在テスト 3-2}\label{tab:mix-3-2}
\begin{tabular}{llcr}
\toprule
\rowcolor{softblue} 分類 & パッケージ & 典型用途 & 負荷 \\
\midrule
数式 & amsmath/mathtools & \makecell{多行整列\\括弧制御} & 高 \\
図 & TikZ/PGFPlots & ベクター図形 & 高 \\
表 & longtable/tabularx & ページ分割表 & 中 \\
コード & listings/fancyvrb & ソース表示 & 中 \\
\bottomrule
\end{tabular}
\end{table}

\subsection{TikZ図}
\begin{figure}[H]
\centering
\begin{tikzpicture}[node distance=18mm,>=Latex]
\node[draw,rounded rectangle,fill=softblue,minimum width=26mm,minimum height=10mm] (src) {Source TeX};
\node[draw,rectangle,fill=softgreen,right=of src,minimum width=26mm,minimum height=10mm] (parse) {Parser};
\node[draw,diamond,aspect=2,fill=softred,right=of parse,minimum width=28mm] (dom) {TeX-DOM};
\node[draw,rounded rectangle,fill=softblue,right=of dom,minimum width=26mm,minimum height=10mm] (pdf) {PDF};
\draw[->,thick] (src) -- node[above]{tokens} (parse);
\draw[->,thick] (parse) -- node[above]{nodes} (dom);
\draw[->,thick] (dom) -- node[above]{layout} (pdf);
\draw[->,bend left=35,dashed] (pdf.north) to node[above]{SyncTeX} (src.north);
\end{tikzpicture}
\caption{TikZによる処理系模式図 3-2}\label{fig:tikz-3-2}
\end{figure}

\subsection{PGFPlots}
\begin{figure}[H]
\centering
\begin{tikzpicture}
\begin{axis}[width=.8\linewidth,height=55mm,grid=major,xlabel={$x$},ylabel={$f(x)$},legend pos=north east]
\addplot[domain=0:6.28,samples=80] {sin(deg(x))};
\addlegendentry{$\sin x$}
\addplot[domain=0:6.28,samples=80,dashed] {cos(deg(x))};
\addlegendentry{$\cos x$}
\addplot[domain=0:6.28,samples=80,dotted] {exp(-x/4)*sin(deg(3*x))};
\addlegendentry{$e^{-x/4}\sin(3x)$}
\end{axis}
\end{tikzpicture}
\caption{PGFPlotsテスト 3-2}\label{fig:pgfplots-3-2}
\end{figure}

\subsection{コードとアルゴリズム}
\begin{lstlisting}[style=pythonstyle,caption={Python風擬似コード 3-2}]
from math import exp

def partition_function(beta, energies):
    return sum(exp(-beta * e) for e in energies)

def free_energy(beta, energies):
    z = partition_function(beta, energies)
    return -log(z) / beta
\end{lstlisting}

\begin{algorithm}[H]
\caption{差分プレビュー更新アルゴリズム 3-2}
\KwIn{old DOM, new source tokens}
\KwOut{minimal changed regions}
Parse changed token interval\;
Update syntax tree nodes\;
Compute layout invalidation frontier\;
\If{math box changed}{Rebuild equation box and update cross references\;}
\Else{Reuse previous node metrics\;}
Return changed PDF rectangles\;
\end{algorithm}

\subsection{化学・回路・箱}
\begin{greenbox}{化学式と単位}
水の生成反応は \ce{2H2 + O2 -> 2H2O} であり、標準状態の議論では温度 \SI{298.15}{\kelvin}、圧力 \SI{1}{\bar} を使う。構造式の簡易表示: \chemfig{H-O-H}。
\end{greenbox}

\begin{figure}[H]
\centering
\begin{circuitikz}[scale=0.9]
\draw (0,0) to[V,l=$V_0$] (0,3) to[R,l=$R$] (3,3) to[C,l=$C$] (3,0) -- (0,0);
\draw (3,3) -- (5,3) to[L,l=$L$] (5,0) -- (3,0);
\end{circuitikz}
\caption{CircuitTikZテスト 3-2}\label{fig:circuit-3-2}
\end{figure}

\begin{boxedtheorem}{局所更新の原理}{local-update-3-2}
完全な再コンパイルではなく、構文木・ボックス・参照関係を局所的に更新できれば、TeX執筆体験はWeb DOMに近づく。
\end{boxedtheorem}

\begin{mdframed}[backgroundcolor=softblue,linecolor=mainblue,roundcorner=6pt]
\blindtext[1]
\end{mdframed}


\section{総合テストセクション 3-3}\label{sec:ch3-s3}
この節では、\TeXDOM、LuaLaTeX、数式組版、図表、相互参照、索引生成を同時に試験する。索引語としてLuaLaTeX\index{LuaLaTeX}、TikZ\index{TikZ}、tcolorbox\index{tcolorbox}、数式\index{数式}、日本語組版\index{日本語組版}を登録する。

\begin{definition}[状態ベクトル]
複素ヒルベルト空間 $\mathcal{H}$ の単位ベクトル $\ket{\psi}\in\mathcal{H}$ を純粋状態という。射影 $\rho=\ket{\psi}\bra{\psi}$ は $\rho^2=\rho$, $\Tr\rho=1$ を満たす。
\end{definition}

\begin{theorem}[スペクトル分解の簡易形]
任意のエルミート行列 $A\in\C^{n\times n}$ には、ユニタリ行列 $U$ と実対角行列 $D=\diag(\lambda_1,\dots,\lambda_n)$ が存在して
\begin{equation}\label{eq:spectral-3-3}
A=UDU^\dagger,\qquad \ee^{-\ii t A}=U\diag(\ee^{-\ii t\lambda_1},\dots,\ee^{-\ii t\lambda_n})U^\dagger
\end{equation}
と書ける。
\end{theorem}

\begin{proof}
有限次元のスペクトル定理による。相異なる固有値に属する固有空間は直交し、各固有空間で正規直交基底を選ぶことで $U$ が得られる。\cref{eq:spectral-3-3} は関数計算の定義から従う。
\end{proof}

\begin{empheq}[box=\fbox]{align}
Z(\beta)&=\Tr \ee^{-\beta H}=\sum_n \ee^{-\beta E_n},\\
F(\beta)&=-\frac{1}{\beta}\log Z(\beta),\\
\langle E\rangle&=-\frac{\partial}{\partial\beta}\log Z(\beta).
\end{empheq}

\begin{bluebox}{日本語・英語・記号混在テスト}
これは長い日本語段落である。The quick brown fox jumps over the lazy dog. 数式 $\int_{-\infty}^{\infty}\ee^{-x^2}\dd x=\sqrt{\pi}$ と、化学式 \ce{2H2 + O2 -> 2H2O} と、単位 \SI{6.62607015e-34}{\joule\second} を同時に含む。さらに、\enquote{引用符}、脚注\footnote{脚注テスト。日本語・English・$E=mc^2$ を含む。}、余白注\marginnote{Margin note 3-3}も入れる。
\end{bluebox}

\subsection{行列・テンソル・微分形式}
\begin{align}
\mat{M}_{3} &=
\begin{pmatrix}
1 & 3 & 3 \\
0 & 1 & 3+3 \\
3 & 0 & 1
\end{pmatrix},
&
\det \mat{M}_{3}&=1-3(3)(3+3),\\
\tensor{R}{^\rho_{\sigma\mu\nu}}&=\partial_\mu \Gamma^\rho_{\nu\sigma}-\partial_\nu\Gamma^\rho_{\mu\sigma}+\Gamma^\rho_{\mu\lambda}\Gamma^\lambda_{\nu\sigma}-\Gamma^\rho_{\nu\lambda}\Gamma^\lambda_{\mu\sigma},\\
\oint_{\partial\Omega}\vect{F}\cdot\vect{n}\dd S&=\iiint_{\Omega}\nabla\cdot\vect{F}\dd V.
\end{align}

\subsection{表と長いセル}
\begin{table}[H]
\centering
\caption{booktabs, makecell, multirow, colortbl の混在テスト 3-3}\label{tab:mix-3-3}
\begin{tabular}{llcr}
\toprule
\rowcolor{softblue} 分類 & パッケージ & 典型用途 & 負荷 \\
\midrule
数式 & amsmath/mathtools & \makecell{多行整列\\括弧制御} & 高 \\
図 & TikZ/PGFPlots & ベクター図形 & 高 \\
表 & longtable/tabularx & ページ分割表 & 中 \\
コード & listings/fancyvrb & ソース表示 & 中 \\
\bottomrule
\end{tabular}
\end{table}

\subsection{TikZ図}
\begin{figure}[H]
\centering
\begin{tikzpicture}[node distance=18mm,>=Latex]
\node[draw,rounded rectangle,fill=softblue,minimum width=26mm,minimum height=10mm] (src) {Source TeX};
\node[draw,rectangle,fill=softgreen,right=of src,minimum width=26mm,minimum height=10mm] (parse) {Parser};
\node[draw,diamond,aspect=2,fill=softred,right=of parse,minimum width=28mm] (dom) {TeX-DOM};
\node[draw,rounded rectangle,fill=softblue,right=of dom,minimum width=26mm,minimum height=10mm] (pdf) {PDF};
\draw[->,thick] (src) -- node[above]{tokens} (parse);
\draw[->,thick] (parse) -- node[above]{nodes} (dom);
\draw[->,thick] (dom) -- node[above]{layout} (pdf);
\draw[->,bend left=35,dashed] (pdf.north) to node[above]{SyncTeX} (src.north);
\end{tikzpicture}
\caption{TikZによる処理系模式図 3-3}\label{fig:tikz-3-3}
\end{figure}

\subsection{PGFPlots}
\begin{figure}[H]
\centering
\begin{tikzpicture}
\begin{axis}[width=.8\linewidth,height=55mm,grid=major,xlabel={$x$},ylabel={$f(x)$},legend pos=north east]
\addplot[domain=0:6.28,samples=80] {sin(deg(x))};
\addlegendentry{$\sin x$}
\addplot[domain=0:6.28,samples=80,dashed] {cos(deg(x))};
\addlegendentry{$\cos x$}
\addplot[domain=0:6.28,samples=80,dotted] {exp(-x/4)*sin(deg(3*x))};
\addlegendentry{$e^{-x/4}\sin(3x)$}
\end{axis}
\end{tikzpicture}
\caption{PGFPlotsテスト 3-3}\label{fig:pgfplots-3-3}
\end{figure}

\subsection{コードとアルゴリズム}
\begin{lstlisting}[style=pythonstyle,caption={Python風擬似コード 3-3}]
from math import exp

def partition_function(beta, energies):
    return sum(exp(-beta * e) for e in energies)

def free_energy(beta, energies):
    z = partition_function(beta, energies)
    return -log(z) / beta
\end{lstlisting}

\begin{algorithm}[H]
\caption{差分プレビュー更新アルゴリズム 3-3}
\KwIn{old DOM, new source tokens}
\KwOut{minimal changed regions}
Parse changed token interval\;
Update syntax tree nodes\;
Compute layout invalidation frontier\;
\If{math box changed}{Rebuild equation box and update cross references\;}
\Else{Reuse previous node metrics\;}
Return changed PDF rectangles\;
\end{algorithm}

\subsection{化学・回路・箱}
\begin{greenbox}{化学式と単位}
水の生成反応は \ce{2H2 + O2 -> 2H2O} であり、標準状態の議論では温度 \SI{298.15}{\kelvin}、圧力 \SI{1}{\bar} を使う。構造式の簡易表示: \chemfig{H-O-H}。
\end{greenbox}

\begin{figure}[H]
\centering
\begin{circuitikz}[scale=0.9]
\draw (0,0) to[V,l=$V_0$] (0,3) to[R,l=$R$] (3,3) to[C,l=$C$] (3,0) -- (0,0);
\draw (3,3) -- (5,3) to[L,l=$L$] (5,0) -- (3,0);
\end{circuitikz}
\caption{CircuitTikZテスト 3-3}\label{fig:circuit-3-3}
\end{figure}

\begin{boxedtheorem}{局所更新の原理}{local-update-3-3}
完全な再コンパイルではなく、構文木・ボックス・参照関係を局所的に更新できれば、TeX執筆体験はWeb DOMに近づく。
\end{boxedtheorem}

\begin{mdframed}[backgroundcolor=softblue,linecolor=mainblue,roundcorner=6pt]
\blindtext[1]
\end{mdframed}

\chapter{巨大総合章 4: パッケージ横断テスト}\label{chap:4}


\begin{redbox}{章の目的}
この章は、同じパターンを少しずつ変えながら反復し、アウトライン、ラベル、図表番号、ページ分割、脚注、索引、長い数式、コード、図形描画を大量に発生させるための章である。
\end{redbox}


\section{総合テストセクション 4-1}\label{sec:ch4-s1}
この節では、\TeXDOM、LuaLaTeX、数式組版、図表、相互参照、索引生成を同時に試験する。索引語としてLuaLaTeX\index{LuaLaTeX}、TikZ\index{TikZ}、tcolorbox\index{tcolorbox}、数式\index{数式}、日本語組版\index{日本語組版}を登録する。

\begin{definition}[状態ベクトル]
複素ヒルベルト空間 $\mathcal{H}$ の単位ベクトル $\ket{\psi}\in\mathcal{H}$ を純粋状態という。射影 $\rho=\ket{\psi}\bra{\psi}$ は $\rho^2=\rho$, $\Tr\rho=1$ を満たす。
\end{definition}

\begin{theorem}[スペクトル分解の簡易形]
任意のエルミート行列 $A\in\C^{n\times n}$ には、ユニタリ行列 $U$ と実対角行列 $D=\diag(\lambda_1,\dots,\lambda_n)$ が存在して
\begin{equation}\label{eq:spectral-4-1}
A=UDU^\dagger,\qquad \ee^{-\ii t A}=U\diag(\ee^{-\ii t\lambda_1},\dots,\ee^{-\ii t\lambda_n})U^\dagger
\end{equation}
と書ける。
\end{theorem}

\begin{proof}
有限次元のスペクトル定理による。相異なる固有値に属する固有空間は直交し、各固有空間で正規直交基底を選ぶことで $U$ が得られる。\cref{eq:spectral-4-1} は関数計算の定義から従う。
\end{proof}

\begin{empheq}[box=\fbox]{align}
Z(\beta)&=\Tr \ee^{-\beta H}=\sum_n \ee^{-\beta E_n},\\
F(\beta)&=-\frac{1}{\beta}\log Z(\beta),\\
\langle E\rangle&=-\frac{\partial}{\partial\beta}\log Z(\beta).
\end{empheq}

\begin{bluebox}{日本語・英語・記号混在テスト}
これは長い日本語段落である。The quick brown fox jumps over the lazy dog. 数式 $\int_{-\infty}^{\infty}\ee^{-x^2}\dd x=\sqrt{\pi}$ と、化学式 \ce{2H2 + O2 -> 2H2O} と、単位 \SI{6.62607015e-34}{\joule\second} を同時に含む。さらに、\enquote{引用符}、脚注\footnote{脚注テスト。日本語・English・$E=mc^2$ を含む。}、余白注\marginnote{Margin note 4-1}も入れる。
\end{bluebox}

\subsection{行列・テンソル・微分形式}
\begin{align}
\mat{M}_{1} &=
\begin{pmatrix}
1 & 1 & 4 \\
0 & 1 & 1+4 \\
4 & 0 & 1
\end{pmatrix},
&
\det \mat{M}_{1}&=1-4(1)(1+4),\\
\tensor{R}{^\rho_{\sigma\mu\nu}}&=\partial_\mu \Gamma^\rho_{\nu\sigma}-\partial_\nu\Gamma^\rho_{\mu\sigma}+\Gamma^\rho_{\mu\lambda}\Gamma^\lambda_{\nu\sigma}-\Gamma^\rho_{\nu\lambda}\Gamma^\lambda_{\mu\sigma},\\
\oint_{\partial\Omega}\vect{F}\cdot\vect{n}\dd S&=\iiint_{\Omega}\nabla\cdot\vect{F}\dd V.
\end{align}

\subsection{表と長いセル}
\begin{table}[H]
\centering
\caption{booktabs, makecell, multirow, colortbl の混在テスト 4-1}\label{tab:mix-4-1}
\begin{tabular}{llcr}
\toprule
\rowcolor{softblue} 分類 & パッケージ & 典型用途 & 負荷 \\
\midrule
数式 & amsmath/mathtools & \makecell{多行整列\\括弧制御} & 高 \\
図 & TikZ/PGFPlots & ベクター図形 & 高 \\
表 & longtable/tabularx & ページ分割表 & 中 \\
コード & listings/fancyvrb & ソース表示 & 中 \\
\bottomrule
\end{tabular}
\end{table}

\subsection{TikZ図}
\begin{figure}[H]
\centering
\begin{tikzpicture}[node distance=18mm,>=Latex]
\node[draw,rounded rectangle,fill=softblue,minimum width=26mm,minimum height=10mm] (src) {Source TeX};
\node[draw,rectangle,fill=softgreen,right=of src,minimum width=26mm,minimum height=10mm] (parse) {Parser};
\node[draw,diamond,aspect=2,fill=softred,right=of parse,minimum width=28mm] (dom) {TeX-DOM};
\node[draw,rounded rectangle,fill=softblue,right=of dom,minimum width=26mm,minimum height=10mm] (pdf) {PDF};
\draw[->,thick] (src) -- node[above]{tokens} (parse);
\draw[->,thick] (parse) -- node[above]{nodes} (dom);
\draw[->,thick] (dom) -- node[above]{layout} (pdf);
\draw[->,bend left=35,dashed] (pdf.north) to node[above]{SyncTeX} (src.north);
\end{tikzpicture}
\caption{TikZによる処理系模式図 4-1}\label{fig:tikz-4-1}
\end{figure}

\subsection{PGFPlots}
\begin{figure}[H]
\centering
\begin{tikzpicture}
\begin{axis}[width=.8\linewidth,height=55mm,grid=major,xlabel={$x$},ylabel={$f(x)$},legend pos=north east]
\addplot[domain=0:6.28,samples=80] {sin(deg(x))};
\addlegendentry{$\sin x$}
\addplot[domain=0:6.28,samples=80,dashed] {cos(deg(x))};
\addlegendentry{$\cos x$}
\addplot[domain=0:6.28,samples=80,dotted] {exp(-x/4)*sin(deg(3*x))};
\addlegendentry{$e^{-x/4}\sin(3x)$}
\end{axis}
\end{tikzpicture}
\caption{PGFPlotsテスト 4-1}\label{fig:pgfplots-4-1}
\end{figure}

\subsection{コードとアルゴリズム}
\begin{lstlisting}[style=pythonstyle,caption={Python風擬似コード 4-1}]
from math import exp

def partition_function(beta, energies):
    return sum(exp(-beta * e) for e in energies)

def free_energy(beta, energies):
    z = partition_function(beta, energies)
    return -log(z) / beta
\end{lstlisting}

\begin{algorithm}[H]
\caption{差分プレビュー更新アルゴリズム 4-1}
\KwIn{old DOM, new source tokens}
\KwOut{minimal changed regions}
Parse changed token interval\;
Update syntax tree nodes\;
Compute layout invalidation frontier\;
\If{math box changed}{Rebuild equation box and update cross references\;}
\Else{Reuse previous node metrics\;}
Return changed PDF rectangles\;
\end{algorithm}

\subsection{化学・回路・箱}
\begin{greenbox}{化学式と単位}
水の生成反応は \ce{2H2 + O2 -> 2H2O} であり、標準状態の議論では温度 \SI{298.15}{\kelvin}、圧力 \SI{1}{\bar} を使う。構造式の簡易表示: \chemfig{H-O-H}。
\end{greenbox}

\begin{figure}[H]
\centering
\begin{circuitikz}[scale=0.9]
\draw (0,0) to[V,l=$V_0$] (0,3) to[R,l=$R$] (3,3) to[C,l=$C$] (3,0) -- (0,0);
\draw (3,3) -- (5,3) to[L,l=$L$] (5,0) -- (3,0);
\end{circuitikz}
\caption{CircuitTikZテスト 4-1}\label{fig:circuit-4-1}
\end{figure}

\begin{boxedtheorem}{局所更新の原理}{local-update-4-1}
完全な再コンパイルではなく、構文木・ボックス・参照関係を局所的に更新できれば、TeX執筆体験はWeb DOMに近づく。
\end{boxedtheorem}

\begin{mdframed}[backgroundcolor=softblue,linecolor=mainblue,roundcorner=6pt]
\blindtext[1]
\end{mdframed}


\section{総合テストセクション 4-2}\label{sec:ch4-s2}
この節では、\TeXDOM、LuaLaTeX、数式組版、図表、相互参照、索引生成を同時に試験する。索引語としてLuaLaTeX\index{LuaLaTeX}、TikZ\index{TikZ}、tcolorbox\index{tcolorbox}、数式\index{数式}、日本語組版\index{日本語組版}を登録する。

\begin{definition}[状態ベクトル]
複素ヒルベルト空間 $\mathcal{H}$ の単位ベクトル $\ket{\psi}\in\mathcal{H}$ を純粋状態という。射影 $\rho=\ket{\psi}\bra{\psi}$ は $\rho^2=\rho$, $\Tr\rho=1$ を満たす。
\end{definition}

\begin{theorem}[スペクトル分解の簡易形]
任意のエルミート行列 $A\in\C^{n\times n}$ には、ユニタリ行列 $U$ と実対角行列 $D=\diag(\lambda_1,\dots,\lambda_n)$ が存在して
\begin{equation}\label{eq:spectral-4-2}
A=UDU^\dagger,\qquad \ee^{-\ii t A}=U\diag(\ee^{-\ii t\lambda_1},\dots,\ee^{-\ii t\lambda_n})U^\dagger
\end{equation}
と書ける。
\end{theorem}

\begin{proof}
有限次元のスペクトル定理による。相異なる固有値に属する固有空間は直交し、各固有空間で正規直交基底を選ぶことで $U$ が得られる。\cref{eq:spectral-4-2} は関数計算の定義から従う。
\end{proof}

\begin{empheq}[box=\fbox]{align}
Z(\beta)&=\Tr \ee^{-\beta H}=\sum_n \ee^{-\beta E_n},\\
F(\beta)&=-\frac{1}{\beta}\log Z(\beta),\\
\langle E\rangle&=-\frac{\partial}{\partial\beta}\log Z(\beta).
\end{empheq}

\begin{bluebox}{日本語・英語・記号混在テスト}
これは長い日本語段落である。The quick brown fox jumps over the lazy dog. 数式 $\int_{-\infty}^{\infty}\ee^{-x^2}\dd x=\sqrt{\pi}$ と、化学式 \ce{2H2 + O2 -> 2H2O} と、単位 \SI{6.62607015e-34}{\joule\second} を同時に含む。さらに、\enquote{引用符}、脚注\footnote{脚注テスト。日本語・English・$E=mc^2$ を含む。}、余白注\marginnote{Margin note 4-2}も入れる。
\end{bluebox}

\subsection{行列・テンソル・微分形式}
\begin{align}
\mat{M}_{2} &=
\begin{pmatrix}
1 & 2 & 4 \\
0 & 1 & 2+4 \\
4 & 0 & 1
\end{pmatrix},
&
\det \mat{M}_{2}&=1-4(2)(2+4),\\
\tensor{R}{^\rho_{\sigma\mu\nu}}&=\partial_\mu \Gamma^\rho_{\nu\sigma}-\partial_\nu\Gamma^\rho_{\mu\sigma}+\Gamma^\rho_{\mu\lambda}\Gamma^\lambda_{\nu\sigma}-\Gamma^\rho_{\nu\lambda}\Gamma^\lambda_{\mu\sigma},\\
\oint_{\partial\Omega}\vect{F}\cdot\vect{n}\dd S&=\iiint_{\Omega}\nabla\cdot\vect{F}\dd V.
\end{align}

\subsection{表と長いセル}
\begin{table}[H]
\centering
\caption{booktabs, makecell, multirow, colortbl の混在テスト 4-2}\label{tab:mix-4-2}
\begin{tabular}{llcr}
\toprule
\rowcolor{softblue} 分類 & パッケージ & 典型用途 & 負荷 \\
\midrule
数式 & amsmath/mathtools & \makecell{多行整列\\括弧制御} & 高 \\
図 & TikZ/PGFPlots & ベクター図形 & 高 \\
表 & longtable/tabularx & ページ分割表 & 中 \\
コード & listings/fancyvrb & ソース表示 & 中 \\
\bottomrule
\end{tabular}
\end{table}

\subsection{TikZ図}
\begin{figure}[H]
\centering
\begin{tikzpicture}[node distance=18mm,>=Latex]
\node[draw,rounded rectangle,fill=softblue,minimum width=26mm,minimum height=10mm] (src) {Source TeX};
\node[draw,rectangle,fill=softgreen,right=of src,minimum width=26mm,minimum height=10mm] (parse) {Parser};
\node[draw,diamond,aspect=2,fill=softred,right=of parse,minimum width=28mm] (dom) {TeX-DOM};
\node[draw,rounded rectangle,fill=softblue,right=of dom,minimum width=26mm,minimum height=10mm] (pdf) {PDF};
\draw[->,thick] (src) -- node[above]{tokens} (parse);
\draw[->,thick] (parse) -- node[above]{nodes} (dom);
\draw[->,thick] (dom) -- node[above]{layout} (pdf);
\draw[->,bend left=35,dashed] (pdf.north) to node[above]{SyncTeX} (src.north);
\end{tikzpicture}
\caption{TikZによる処理系模式図 4-2}\label{fig:tikz-4-2}
\end{figure}

\subsection{PGFPlots}
\begin{figure}[H]
\centering
\begin{tikzpicture}
\begin{axis}[width=.8\linewidth,height=55mm,grid=major,xlabel={$x$},ylabel={$f(x)$},legend pos=north east]
\addplot[domain=0:6.28,samples=80] {sin(deg(x))};
\addlegendentry{$\sin x$}
\addplot[domain=0:6.28,samples=80,dashed] {cos(deg(x))};
\addlegendentry{$\cos x$}
\addplot[domain=0:6.28,samples=80,dotted] {exp(-x/4)*sin(deg(3*x))};
\addlegendentry{$e^{-x/4}\sin(3x)$}
\end{axis}
\end{tikzpicture}
\caption{PGFPlotsテスト 4-2}\label{fig:pgfplots-4-2}
\end{figure}

\subsection{コードとアルゴリズム}
\begin{lstlisting}[style=pythonstyle,caption={Python風擬似コード 4-2}]
from math import exp

def partition_function(beta, energies):
    return sum(exp(-beta * e) for e in energies)

def free_energy(beta, energies):
    z = partition_function(beta, energies)
    return -log(z) / beta
\end{lstlisting}

\begin{algorithm}[H]
\caption{差分プレビュー更新アルゴリズム 4-2}
\KwIn{old DOM, new source tokens}
\KwOut{minimal changed regions}
Parse changed token interval\;
Update syntax tree nodes\;
Compute layout invalidation frontier\;
\If{math box changed}{Rebuild equation box and update cross references\;}
\Else{Reuse previous node metrics\;}
Return changed PDF rectangles\;
\end{algorithm}

\subsection{化学・回路・箱}
\begin{greenbox}{化学式と単位}
水の生成反応は \ce{2H2 + O2 -> 2H2O} であり、標準状態の議論では温度 \SI{298.15}{\kelvin}、圧力 \SI{1}{\bar} を使う。構造式の簡易表示: \chemfig{H-O-H}。
\end{greenbox}

\begin{figure}[H]
\centering
\begin{circuitikz}[scale=0.9]
\draw (0,0) to[V,l=$V_0$] (0,3) to[R,l=$R$] (3,3) to[C,l=$C$] (3,0) -- (0,0);
\draw (3,3) -- (5,3) to[L,l=$L$] (5,0) -- (3,0);
\end{circuitikz}
\caption{CircuitTikZテスト 4-2}\label{fig:circuit-4-2}
\end{figure}

\begin{boxedtheorem}{局所更新の原理}{local-update-4-2}
完全な再コンパイルではなく、構文木・ボックス・参照関係を局所的に更新できれば、TeX執筆体験はWeb DOMに近づく。
\end{boxedtheorem}

\begin{mdframed}[backgroundcolor=softblue,linecolor=mainblue,roundcorner=6pt]
\blindtext[1]
\end{mdframed}


\section{総合テストセクション 4-3}\label{sec:ch4-s3}
この節では、\TeXDOM、LuaLaTeX、数式組版、図表、相互参照、索引生成を同時に試験する。索引語としてLuaLaTeX\index{LuaLaTeX}、TikZ\index{TikZ}、tcolorbox\index{tcolorbox}、数式\index{数式}、日本語組版\index{日本語組版}を登録する。

\begin{definition}[状態ベクトル]
複素ヒルベルト空間 $\mathcal{H}$ の単位ベクトル $\ket{\psi}\in\mathcal{H}$ を純粋状態という。射影 $\rho=\ket{\psi}\bra{\psi}$ は $\rho^2=\rho$, $\Tr\rho=1$ を満たす。
\end{definition}

\begin{theorem}[スペクトル分解の簡易形]
任意のエルミート行列 $A\in\C^{n\times n}$ には、ユニタリ行列 $U$ と実対角行列 $D=\diag(\lambda_1,\dots,\lambda_n)$ が存在して
\begin{equation}\label{eq:spectral-4-3}
A=UDU^\dagger,\qquad \ee^{-\ii t A}=U\diag(\ee^{-\ii t\lambda_1},\dots,\ee^{-\ii t\lambda_n})U^\dagger
\end{equation}
と書ける。
\end{theorem}

\begin{proof}
有限次元のスペクトル定理による。相異なる固有値に属する固有空間は直交し、各固有空間で正規直交基底を選ぶことで $U$ が得られる。\cref{eq:spectral-4-3} は関数計算の定義から従う。
\end{proof}

\begin{empheq}[box=\fbox]{align}
Z(\beta)&=\Tr \ee^{-\beta H}=\sum_n \ee^{-\beta E_n},\\
F(\beta)&=-\frac{1}{\beta}\log Z(\beta),\\
\langle E\rangle&=-\frac{\partial}{\partial\beta}\log Z(\beta).
\end{empheq}

\begin{bluebox}{日本語・英語・記号混在テスト}
これは長い日本語段落である。The quick brown fox jumps over the lazy dog. 数式 $\int_{-\infty}^{\infty}\ee^{-x^2}\dd x=\sqrt{\pi}$ と、化学式 \ce{2H2 + O2 -> 2H2O} と、単位 \SI{6.62607015e-34}{\joule\second} を同時に含む。さらに、\enquote{引用符}、脚注\footnote{脚注テスト。日本語・English・$E=mc^2$ を含む。}、余白注\marginnote{Margin note 4-3}も入れる。
\end{bluebox}

\subsection{行列・テンソル・微分形式}
\begin{align}
\mat{M}_{3} &=
\begin{pmatrix}
1 & 3 & 4 \\
0 & 1 & 3+4 \\
4 & 0 & 1
\end{pmatrix},
&
\det \mat{M}_{3}&=1-4(3)(3+4),\\
\tensor{R}{^\rho_{\sigma\mu\nu}}&=\partial_\mu \Gamma^\rho_{\nu\sigma}-\partial_\nu\Gamma^\rho_{\mu\sigma}+\Gamma^\rho_{\mu\lambda}\Gamma^\lambda_{\nu\sigma}-\Gamma^\rho_{\nu\lambda}\Gamma^\lambda_{\mu\sigma},\\
\oint_{\partial\Omega}\vect{F}\cdot\vect{n}\dd S&=\iiint_{\Omega}\nabla\cdot\vect{F}\dd V.
\end{align}

\subsection{表と長いセル}
\begin{table}[H]
\centering
\caption{booktabs, makecell, multirow, colortbl の混在テスト 4-3}\label{tab:mix-4-3}
\begin{tabular}{llcr}
\toprule
\rowcolor{softblue} 分類 & パッケージ & 典型用途 & 負荷 \\
\midrule
数式 & amsmath/mathtools & \makecell{多行整列\\括弧制御} & 高 \\
図 & TikZ/PGFPlots & ベクター図形 & 高 \\
表 & longtable/tabularx & ページ分割表 & 中 \\
コード & listings/fancyvrb & ソース表示 & 中 \\
\bottomrule
\end{tabular}
\end{table}

\subsection{TikZ図}
\begin{figure}[H]
\centering
\begin{tikzpicture}[node distance=18mm,>=Latex]
\node[draw,rounded rectangle,fill=softblue,minimum width=26mm,minimum height=10mm] (src) {Source TeX};
\node[draw,rectangle,fill=softgreen,right=of src,minimum width=26mm,minimum height=10mm] (parse) {Parser};
\node[draw,diamond,aspect=2,fill=softred,right=of parse,minimum width=28mm] (dom) {TeX-DOM};
\node[draw,rounded rectangle,fill=softblue,right=of dom,minimum width=26mm,minimum height=10mm] (pdf) {PDF};
\draw[->,thick] (src) -- node[above]{tokens} (parse);
\draw[->,thick] (parse) -- node[above]{nodes} (dom);
\draw[->,thick] (dom) -- node[above]{layout} (pdf);
\draw[->,bend left=35,dashed] (pdf.north) to node[above]{SyncTeX} (src.north);
\end{tikzpicture}
\caption{TikZによる処理系模式図 4-3}\label{fig:tikz-4-3}
\end{figure}

\subsection{PGFPlots}
\begin{figure}[H]
\centering
\begin{tikzpicture}
\begin{axis}[width=.8\linewidth,height=55mm,grid=major,xlabel={$x$},ylabel={$f(x)$},legend pos=north east]
\addplot[domain=0:6.28,samples=80] {sin(deg(x))};
\addlegendentry{$\sin x$}
\addplot[domain=0:6.28,samples=80,dashed] {cos(deg(x))};
\addlegendentry{$\cos x$}
\addplot[domain=0:6.28,samples=80,dotted] {exp(-x/4)*sin(deg(3*x))};
\addlegendentry{$e^{-x/4}\sin(3x)$}
\end{axis}
\end{tikzpicture}
\caption{PGFPlotsテスト 4-3}\label{fig:pgfplots-4-3}
\end{figure}

\subsection{コードとアルゴリズム}
\begin{lstlisting}[style=pythonstyle,caption={Python風擬似コード 4-3}]
from math import exp

def partition_function(beta, energies):
    return sum(exp(-beta * e) for e in energies)

def free_energy(beta, energies):
    z = partition_function(beta, energies)
    return -log(z) / beta
\end{lstlisting}

\begin{algorithm}[H]
\caption{差分プレビュー更新アルゴリズム 4-3}
\KwIn{old DOM, new source tokens}
\KwOut{minimal changed regions}
Parse changed token interval\;
Update syntax tree nodes\;
Compute layout invalidation frontier\;
\If{math box changed}{Rebuild equation box and update cross references\;}
\Else{Reuse previous node metrics\;}
Return changed PDF rectangles\;
\end{algorithm}

\subsection{化学・回路・箱}
\begin{greenbox}{化学式と単位}
水の生成反応は \ce{2H2 + O2 -> 2H2O} であり、標準状態の議論では温度 \SI{298.15}{\kelvin}、圧力 \SI{1}{\bar} を使う。構造式の簡易表示: \chemfig{H-O-H}。
\end{greenbox}

\begin{figure}[H]
\centering
\begin{circuitikz}[scale=0.9]
\draw (0,0) to[V,l=$V_0$] (0,3) to[R,l=$R$] (3,3) to[C,l=$C$] (3,0) -- (0,0);
\draw (3,3) -- (5,3) to[L,l=$L$] (5,0) -- (3,0);
\end{circuitikz}
\caption{CircuitTikZテスト 4-3}\label{fig:circuit-4-3}
\end{figure}

\begin{boxedtheorem}{局所更新の原理}{local-update-4-3}
完全な再コンパイルではなく、構文木・ボックス・参照関係を局所的に更新できれば、TeX執筆体験はWeb DOMに近づく。
\end{boxedtheorem}

\begin{mdframed}[backgroundcolor=softblue,linecolor=mainblue,roundcorner=6pt]
\blindtext[1]
\end{mdframed}

\appendix
\chapter{横向きページと長表}
\begin{landscape}
\begin{longtable}{llp{55mm}p{80mm}r}
\caption{非常に長いパッケージ用途一覧}\label{tab:long-package-list}\\
\toprule
番号 & 分類 & パッケージ & 主な用途 & 優先度\\
\midrule
\endfirsthead
\toprule
番号 & 分類 & パッケージ & 主な用途 & 優先度\\
\midrule
\endhead
1 & 日本語 & \texttt{luatexja} & LuaLaTeXで日本語を自然に組む。巨大文書では、相互作用・警告・ページ分割・再コンパイル時の挙動を確認する。 & 2\\
2 & フォント & \texttt{fontspec} & OpenTypeフォント利用。巨大文書では、相互作用・警告・ページ分割・再コンパイル時の挙動を確認する。 & 3\\
3 & 数式 & \texttt{amsmath} & 整列数式。巨大文書では、相互作用・警告・ページ分割・再コンパイル時の挙動を確認する。 & 4\\
4 & 数式 & \texttt{mathtools} & amsmath拡張。巨大文書では、相互作用・警告・ページ分割・再コンパイル時の挙動を確認する。 & 5\\
5 & 定理 & \texttt{amsthm} & 定理環境。巨大文書では、相互作用・警告・ページ分割・再コンパイル時の挙動を確認する。 & 1\\
6 & 図 & \texttt{tikz} & ベクター作図。巨大文書では、相互作用・警告・ページ分割・再コンパイル時の挙動を確認する。 & 2\\
7 & グラフ & \texttt{pgfplots} & 関数グラフ。巨大文書では、相互作用・警告・ページ分割・再コンパイル時の挙動を確認する。 & 3\\
8 & 表 & \texttt{booktabs} & 美しい罫線。巨大文書では、相互作用・警告・ページ分割・再コンパイル時の挙動を確認する。 & 4\\
9 & 表 & \texttt{longtable} & ページ分割表。巨大文書では、相互作用・警告・ページ分割・再コンパイル時の挙動を確認する。 & 5\\
10 & 表 & \texttt{tabularx} & 幅自動調整。巨大文書では、相互作用・警告・ページ分割・再コンパイル時の挙動を確認する。 & 1\\
11 & 箱 & \texttt{tcolorbox} & 高機能ボックス。巨大文書では、相互作用・警告・ページ分割・再コンパイル時の挙動を確認する。 & 2\\
12 & コード & \texttt{listings} & コード表示。巨大文書では、相互作用・警告・ページ分割・再コンパイル時の挙動を確認する。 & 3\\
13 & 疑似コード & \texttt{algorithm2e} & アルゴリズム表示。巨大文書では、相互作用・警告・ページ分割・再コンパイル時の挙動を確認する。 & 4\\
14 & 単位 & \texttt{siunitx} & 単位と数値。巨大文書では、相互作用・警告・ページ分割・再コンパイル時の挙動を確認する。 & 5\\
15 & 化学 & \texttt{mhchem} & 化学式。巨大文書では、相互作用・警告・ページ分割・再コンパイル時の挙動を確認する。 & 1\\
16 & 回路 & \texttt{circuitikz} & 電気回路。巨大文書では、相互作用・警告・ページ分割・再コンパイル時の挙動を確認する。 & 2\\
17 & 参照 & \texttt{hyperref} & リンク。巨大文書では、相互作用・警告・ページ分割・再コンパイル時の挙動を確認する。 & 3\\
18 & 参照 & \texttt{cleveref} & 賢い相互参照。巨大文書では、相互作用・警告・ページ分割・再コンパイル時の挙動を確認する。 & 4\\
19 & 日本語 & \texttt{luatexja} & LuaLaTeXで日本語を自然に組む。巨大文書では、相互作用・警告・ページ分割・再コンパイル時の挙動を確認する。 & 5\\
20 & フォント & \texttt{fontspec} & OpenTypeフォント利用。巨大文書では、相互作用・警告・ページ分割・再コンパイル時の挙動を確認する。 & 1\\
21 & 数式 & \texttt{amsmath} & 整列数式。巨大文書では、相互作用・警告・ページ分割・再コンパイル時の挙動を確認する。 & 2\\
22 & 数式 & \texttt{mathtools} & amsmath拡張。巨大文書では、相互作用・警告・ページ分割・再コンパイル時の挙動を確認する。 & 3\\
23 & 定理 & \texttt{amsthm} & 定理環境。巨大文書では、相互作用・警告・ページ分割・再コンパイル時の挙動を確認する。 & 4\\
24 & 図 & \texttt{tikz} & ベクター作図。巨大文書では、相互作用・警告・ページ分割・再コンパイル時の挙動を確認する。 & 5\\
25 & グラフ & \texttt{pgfplots} & 関数グラフ。巨大文書では、相互作用・警告・ページ分割・再コンパイル時の挙動を確認する。 & 1\\
26 & 表 & \texttt{booktabs} & 美しい罫線。巨大文書では、相互作用・警告・ページ分割・再コンパイル時の挙動を確認する。 & 2\\
27 & 表 & \texttt{longtable} & ページ分割表。巨大文書では、相互作用・警告・ページ分割・再コンパイル時の挙動を確認する。 & 3\\
28 & 表 & \texttt{tabularx} & 幅自動調整。巨大文書では、相互作用・警告・ページ分割・再コンパイル時の挙動を確認する。 & 4\\
29 & 箱 & \texttt{tcolorbox} & 高機能ボックス。巨大文書では、相互作用・警告・ページ分割・再コンパイル時の挙動を確認する。 & 5\\
30 & コード & \texttt{listings} & コード表示。巨大文書では、相互作用・警告・ページ分割・再コンパイル時の挙動を確認する。 & 1\\
31 & 疑似コード & \texttt{algorithm2e} & アルゴリズム表示。巨大文書では、相互作用・警告・ページ分割・再コンパイル時の挙動を確認する。 & 2\\
32 & 単位 & \texttt{siunitx} & 単位と数値。巨大文書では、相互作用・警告・ページ分割・再コンパイル時の挙動を確認する。 & 3\\
33 & 化学 & \texttt{mhchem} & 化学式。巨大文書では、相互作用・警告・ページ分割・再コンパイル時の挙動を確認する。 & 4\\
34 & 回路 & \texttt{circuitikz} & 電気回路。巨大文書では、相互作用・警告・ページ分割・再コンパイル時の挙動を確認する。 & 5\\
35 & 参照 & \texttt{hyperref} & リンク。巨大文書では、相互作用・警告・ページ分割・再コンパイル時の挙動を確認する。 & 1\\
36 & 参照 & \texttt{cleveref} & 賢い相互参照。巨大文書では、相互作用・警告・ページ分割・再コンパイル時の挙動を確認する。 & 2\\
37 & 日本語 & \texttt{luatexja} & LuaLaTeXで日本語を自然に組む。巨大文書では、相互作用・警告・ページ分割・再コンパイル時の挙動を確認する。 & 3\\
38 & フォント & \texttt{fontspec} & OpenTypeフォント利用。巨大文書では、相互作用・警告・ページ分割・再コンパイル時の挙動を確認する。 & 4\\
39 & 数式 & \texttt{amsmath} & 整列数式。巨大文書では、相互作用・警告・ページ分割・再コンパイル時の挙動を確認する。 & 5\\
40 & 数式 & \texttt{mathtools} & amsmath拡張。巨大文書では、相互作用・警告・ページ分割・再コンパイル時の挙動を確認する。 & 1\\
41 & 定理 & \texttt{amsthm} & 定理環境。巨大文書では、相互作用・警告・ページ分割・再コンパイル時の挙動を確認する。 & 2\\
42 & 図 & \texttt{tikz} & ベクター作図。巨大文書では、相互作用・警告・ページ分割・再コンパイル時の挙動を確認する。 & 3\\
43 & グラフ & \texttt{pgfplots} & 関数グラフ。巨大文書では、相互作用・警告・ページ分割・再コンパイル時の挙動を確認する。 & 4\\
44 & 表 & \texttt{booktabs} & 美しい罫線。巨大文書では、相互作用・警告・ページ分割・再コンパイル時の挙動を確認する。 & 5\\
45 & 表 & \texttt{longtable} & ページ分割表。巨大文書では、相互作用・警告・ページ分割・再コンパイル時の挙動を確認する。 & 1\\
46 & 表 & \texttt{tabularx} & 幅自動調整。巨大文書では、相互作用・警告・ページ分割・再コンパイル時の挙動を確認する。 & 2\\
47 & 箱 & \texttt{tcolorbox} & 高機能ボックス。巨大文書では、相互作用・警告・ページ分割・再コンパイル時の挙動を確認する。 & 3\\
48 & コード & \texttt{listings} & コード表示。巨大文書では、相互作用・警告・ページ分割・再コンパイル時の挙動を確認する。 & 4\\
49 & 疑似コード & \texttt{algorithm2e} & アルゴリズム表示。巨大文書では、相互作用・警告・ページ分割・再コンパイル時の挙動を確認する。 & 5\\
50 & 単位 & \texttt{siunitx} & 単位と数値。巨大文書では、相互作用・警告・ページ分割・再コンパイル時の挙動を確認する。 & 1\\
51 & 化学 & \texttt{mhchem} & 化学式。巨大文書では、相互作用・警告・ページ分割・再コンパイル時の挙動を確認する。 & 2\\
52 & 回路 & \texttt{circuitikz} & 電気回路。巨大文書では、相互作用・警告・ページ分割・再コンパイル時の挙動を確認する。 & 3\\
53 & 参照 & \texttt{hyperref} & リンク。巨大文書では、相互作用・警告・ページ分割・再コンパイル時の挙動を確認する。 & 4\\
54 & 参照 & \texttt{cleveref} & 賢い相互参照。巨大文書では、相互作用・警告・ページ分割・再コンパイル時の挙動を確認する。 & 5\\
55 & 日本語 & \texttt{luatexja} & LuaLaTeXで日本語を自然に組む。巨大文書では、相互作用・警告・ページ分割・再コンパイル時の挙動を確認する。 & 1\\
56 & フォント & \texttt{fontspec} & OpenTypeフォント利用。巨大文書では、相互作用・警告・ページ分割・再コンパイル時の挙動を確認する。 & 2\\
57 & 数式 & \texttt{amsmath} & 整列数式。巨大文書では、相互作用・警告・ページ分割・再コンパイル時の挙動を確認する。 & 3\\
58 & 数式 & \texttt{mathtools} & amsmath拡張。巨大文書では、相互作用・警告・ページ分割・再コンパイル時の挙動を確認する。 & 4\\
59 & 定理 & \texttt{amsthm} & 定理環境。巨大文書では、相互作用・警告・ページ分割・再コンパイル時の挙動を確認する。 & 5\\
60 & 図 & \texttt{tikz} & ベクター作図。巨大文書では、相互作用・警告・ページ分割・再コンパイル時の挙動を確認する。 & 1\\
61 & グラフ & \texttt{pgfplots} & 関数グラフ。巨大文書では、相互作用・警告・ページ分割・再コンパイル時の挙動を確認する。 & 2\\
62 & 表 & \texttt{booktabs} & 美しい罫線。巨大文書では、相互作用・警告・ページ分割・再コンパイル時の挙動を確認する。 & 3\\
63 & 表 & \texttt{longtable} & ページ分割表。巨大文書では、相互作用・警告・ページ分割・再コンパイル時の挙動を確認する。 & 4\\
64 & 表 & \texttt{tabularx} & 幅自動調整。巨大文書では、相互作用・警告・ページ分割・再コンパイル時の挙動を確認する。 & 5\\
65 & 箱 & \texttt{tcolorbox} & 高機能ボックス。巨大文書では、相互作用・警告・ページ分割・再コンパイル時の挙動を確認する。 & 1\\
66 & コード & \texttt{listings} & コード表示。巨大文書では、相互作用・警告・ページ分割・再コンパイル時の挙動を確認する。 & 2\\
67 & 疑似コード & \texttt{algorithm2e} & アルゴリズム表示。巨大文書では、相互作用・警告・ページ分割・再コンパイル時の挙動を確認する。 & 3\\
68 & 単位 & \texttt{siunitx} & 単位と数値。巨大文書では、相互作用・警告・ページ分割・再コンパイル時の挙動を確認する。 & 4\\
69 & 化学 & \texttt{mhchem} & 化学式。巨大文書では、相互作用・警告・ページ分割・再コンパイル時の挙動を確認する。 & 5\\
70 & 回路 & \texttt{circuitikz} & 電気回路。巨大文書では、相互作用・警告・ページ分割・再コンパイル時の挙動を確認する。 & 1\\
71 & 参照 & \texttt{hyperref} & リンク。巨大文書では、相互作用・警告・ページ分割・再コンパイル時の挙動を確認する。 & 2\\
72 & 参照 & \texttt{cleveref} & 賢い相互参照。巨大文書では、相互作用・警告・ページ分割・再コンパイル時の挙動を確認する。 & 3\\
73 & 日本語 & \texttt{luatexja} & LuaLaTeXで日本語を自然に組む。巨大文書では、相互作用・警告・ページ分割・再コンパイル時の挙動を確認する。 & 4\\
74 & フォント & \texttt{fontspec} & OpenTypeフォント利用。巨大文書では、相互作用・警告・ページ分割・再コンパイル時の挙動を確認する。 & 5\\
75 & 数式 & \texttt{amsmath} & 整列数式。巨大文書では、相互作用・警告・ページ分割・再コンパイル時の挙動を確認する。 & 1\\
76 & 数式 & \texttt{mathtools} & amsmath拡張。巨大文書では、相互作用・警告・ページ分割・再コンパイル時の挙動を確認する。 & 2\\
77 & 定理 & \texttt{amsthm} & 定理環境。巨大文書では、相互作用・警告・ページ分割・再コンパイル時の挙動を確認する。 & 3\\
78 & 図 & \texttt{tikz} & ベクター作図。巨大文書では、相互作用・警告・ページ分割・再コンパイル時の挙動を確認する。 & 4\\
79 & グラフ & \texttt{pgfplots} & 関数グラフ。巨大文書では、相互作用・警告・ページ分割・再コンパイル時の挙動を確認する。 & 5\\
80 & 表 & \texttt{booktabs} & 美しい罫線。巨大文書では、相互作用・警告・ページ分割・再コンパイル時の挙動を確認する。 & 1\\
\bottomrule
\end{longtable}
\end{landscape}

\chapter{二段組・並列段組・引用}
\begin{multicols}{2}
\blindtext[6]
\end{multicols}

\begin{paracol}{2}
\switchcolumn[0]*
左カラム。日本語の長文、数式 $\sum_{n=1}^{\infty}1/n^2=\pi^2/6$、脚注、索引を含む。\index{並列段組}
\switchcolumn
Right column. This column checks paracol with English text, inline math $\nabla\times\vect{E}=-\partial_t\vect{B}$, and references to \cref{tab:long-package-list}.
\switchcolumn[0]
さらに左。\blindtext[1]
\switchcolumn
Further right. \blindtext[1]
\end{paracol}

\chapter{参考文献風リストと索引}
本文中で参照した概念をまとめる。これはBibTeX/Biberなしで単体コンパイルできるように、thebibliography環境で書いている。
\begin{thebibliography}{99}
\bibitem{knuth1984} Donald E. Knuth, \textit{The TeXbook}, Addison-Wesley, 1984.
\bibitem{latexcompanion} Frank Mittelbach and Michel Goossens, \textit{The LaTeX Companion}, Addison-Wesley.
\bibitem{tikzpgf} Till Tantau, \textit{The TikZ and PGF Packages Manual}.
\bibitem{luatexja} LuaTeX-ja project, LuaTeX-ja package documentation.
\end{thebibliography}

\backmatter
\printindex
\end{document}
